\documentclass[conference]{IEEEtran}

\usepackage{fontspec}
\setmainfont{TeX Gyre Termes}
\setsansfont{TeX Gyre Heros}
\setmonofont{DejaVu Sans Mono}[Scale=MatchLowercase]
\usepackage{amsmath,amssymb,amsthm,mathtools}
\usepackage{graphicx}
\usepackage{tikz}
\usetikzlibrary{arrows.meta,positioning}
\usepackage{booktabs}
\usepackage{array}
\usepackage{mathpartir}
\usepackage{mathrsfs}
\usepackage{stmaryrd}
\usepackage{xcolor}
\usepackage{url}
\usepackage{algorithm}
\usepackage{algpseudocode}
\usepackage{placeins}
\usepackage{minted}
\usepackage{hyperref}
\usepackage[capitalise,nameinlink,noabbrev]{cleveref}

\hypersetup{hidelinks}
\urlstyle{same}
\makeatletter
\providecommand{\theHALG@line}{}
\renewcommand{\theHALG@line}{\thealgorithm.\arabic{ALG@line}}
\makeatother

% Keep code excerpts compact enough for the IEEE two-column appendix.
% Reserve space inside the column for line numbers; FancyVerb otherwise puts
% them in the outer margin, where right-column listings overlap the left column.
\setminted{
  fontsize=\scriptsize,
  breaklines=true,
  breakanywhere=true,
  autogobble=true,
  xleftmargin=2em,
  numbersep=0.5em
}


\title{Verified Pythagorean Composition for Adaptive Cryptographic Games: Noise Flooding in Approximate Homomorphic Encryption}

\author{\IEEEauthorblockN{Anonymous submission}}

%\newcommand{\Draft}{}
%\newcommand{\CameraReady}{}

\ifdefined\Draft
\newcommand{\Ethan}[1]{{\color{cyan}[Yi: #1]}}
\newcommand{\Alex}[1]{{\color{magenta}[Alex: #1]}}
\newcommand{\Codex}[1]{{\color{blue}[Codex: #1]}}
\newcommand{\Junyi}[1]{{\color{red}[Junyi: #1]}}
\else
\newcommand{\Ethan}[1]{}
\newcommand{\Alex}[1]{}
\newcommand{\Codex}[1]{}
\newcommand{\Junyi}[1]{}
\fi


\theoremstyle{definition}
\theoremstyle{definition}

\newtheorem{definition}{Definition}[section]

\newaliascnt{theorem}{definition}
\newtheorem{theorem}[theorem]{Theorem}
\aliascntresetthe{theorem}

\newaliascnt{lemma}{definition}
\newtheorem{lemma}[lemma]{Lemma}
\aliascntresetthe{lemma}

\newaliascnt{corollary}{definition}
\newtheorem{corollary}[corollary]{Corollary}
\aliascntresetthe{corollary}

\newcommand{\cA}{\mathcal{A}}
\newcommand{\cB}{\mathcal{B}}
\newcommand{\cD}{\mathcal{D}}
\newcommand{\cP}{\mathcal{P}}
\newcommand{\cQ}{\mathcal{Q}}
\newcommand{\cM}{\mathcal{M}}
\newcommand{\cI}{\mathcal{I}}
\newcommand{\bbR}{\mathbb{R}}
\newcommand{\bbN}{\mathbb{N}}
\newcommand{\bbZ}{\mathbb{Z}}
\newcommand{\parens}[1]{\left(#1\right)}
\newcommand{\norm}[1]{\left\lVert#1\right\rVert}
\newcommand{\dKL}{D_{\mathrm{KL}}}
\newcommand{\KL}[2]{\dKL\!\parens{#1\,\|\,#2}}
\newcommand{\DG}[2]{D_{\bbZ,#1,#2}}
\newcommand{\distr}{\mathsf{Distr}}
\newcommand{\semantic}[1]{\llbracket #1\rrbracket}
\newcommand{\csemantic}[1]{\llbracket #1\rrbracket^{\bot}}
\newcommand{\completed}[1]{\mathsf{Comp}\parens{#1}}
\newcommand{\supp}{\mathsf{supp}}
\newcommand{\TVD}{\Delta}
\newcommand{\Pyth}{\mathsf{Pyth}}
\newcommand{\AEJ}{\mathsf{AE}}
\newcommand{\WinProb}{\mathsf{WinProb}}
\newcommand{\apRHL}{\mathsf{apRHL}}
\newcommand{\KLs}{\mathsf{KLs}}
\newcommand{\code}{\mathsf{code}}
\newcommand{\kvalid}{\mathsf{keys\_valid}}
\newcommand{\tvalid}{\mathsf{table\_valid}}
\newcommand{\doubleplus}{\mathbin{+\!\!+}}
\DeclareMathOperator{\Bern}{Bern}


\begin{document}

\maketitle

\begin{abstract}
Noise flooding is a standard defense against decryption attacks on approximate
homomorphic encryption, but its security proof is unusually sensitive to
composition.
Replacing each of \(q\) adaptive decryption answers with a
statistically close simulation and applying an ordinary hybrid argument loses
linearly in \(q\).
The cryptographic proof instead accumulates conditional
Kullback--Leibler (KL) costs and converts to statistical distance once, giving
the parameter-critical square-root loss.

We machine-checked this argument using Rocq and SSProve.
Given any fully homomorphic encryption scheme that is approximately correct and IND-CPA secure,
we formalize a reduction for every \(q\)-query IND-CPAD adversary and prove
\[
 \Pr[\mathsf{IND\text{-}CPAD}_{\mathsf{NF}}^{\mathcal A}=1]
 \leq \beta_{\mathsf{CPA}}(\mathcal B_{\mathcal A,q})
       + \frac{\sqrt{qn}}{2\gamma}.
\]
where \(n\) is the plaintext dimension and \(\gamma\) is the flooding-width
multiplier.
Our proof constructs a new relational program logic
over SSProve semantics.
Its Pythagorean judgment composes
conditional KL budgets without converting them to statistical distance, and a
verified trace compiler lifts a local oracle rule to arbitrary adaptive
programs with a single final conversion.
\iffalse
The development also supplies the
required discrete-Gaussian analysis.
The checked theorem
isolates the ``good key pair" game hop: imperfect correctness can later be captured
using an additional standard up-to-bad game hop.
As in SSProve, polynomial-time preservation remains the conventional
metatheoretic audit outside the program semantics.
\fi
\end{abstract}

\begin{IEEEkeywords}
approximate homomorphic encryption, noise flooding, IND-CPAD, SSProve,
Rocq, relational program logic, game-based cryptography,
Kullback--Leibler divergence
\end{IEEEkeywords}

\Ethan{Original intro here, Alex's version afterwards}

\section{Introduction}
\label{sec:introduction}

Approximate homomorphic encryption (FHE) makes encrypted real-number computation practical by allowing controlled approximation error.
That error is also relevant in the context of security analysis.
Li and Micciancio showed that observing suitably chosen approximate decryptions can reveal the secret key of CKKS-like schemes~\cite{lm,ckks}; related attacks now reach aggressively parameterized exact FHE as well~\cite{cheon2024attacks}.
In approximate FHE one cannot simply require every decrypted result to equal an exact message: approximation is the advertised interface.

Li, Micciancio, Schultz-Wu, and Sorrell (LMSS) formalized this attack vector through the \emph{IND-CPAD} security game, and analyzed a noise-flooding defense~\cite{lmss}.
The defense perturbs each admissible decryption answer with a fresh, sufficiently wide discrete Gaussian.
On the other hand, the decryption result is useful only when that Gaussian width is small enough to preserve precision, so this parameter choice has itself become a research topic \cite{costache2023precision, bergamaschi2025revisiting}.

The parameter trade-off also exposes a mechanization problem.
Suppose a one-query replacement costs \(\delta\) in statistical distance, then a standard game hop argument (i.e. using data processing inequality) would pay \(q\delta\).
The LMSS analysis instead bounds a conditional KL divergence at each query and applies the Micciancio--Walter Pythagorean preservation argument only after constructing the whole transcript~\cite{mw17}, paying \(O(\sqrt q)\) in statistical distance instead.
For fixed target advantage, this is the difference between choosing flooding width proportional to \(\sqrt q\) and to \(q\), with a corresponding impact on precision and FHE parameters.

We verify that security argument in Rocq on top of SSProve, a framework for machine-checked, game-based cryptography~\cite{ssprove}.
The top-level theorem has ordinary cryptographic semantics: it bounds the winning probability of a concrete IND-CPAD game by the IND-CPA bound of a constructed reduction plus an explicit loss.
Along the way, we construct a new relational program logic and verify its soundness, and implement a verified compiler that transforms arbitrary SSProve code into those compatible with our program logic.

\subsection{Result and security significance}

Let \(S\) be an approximate-FHE scheme whose message space has local charts in
\(\bbZ^n\), let \(\mathsf{NF}_{\gamma}[S]\) be its noise-flooded transform, and
let \(\cA\) make at most \(q\) non-failing decryption queries in the IND-CPAD
game.
Under the explicit assumptions summarized in
\Cref{tab:theorem-boundary}, the formal development constructs an IND-CPA
adversary \(\cB_{\cA,q}\) and proves
\begin{equation}
\label{eq:headline}
 \WinProb[\mathsf{IND\text{-}CPAD}_{\mathsf{NF}_{\gamma}[S]}(\cA)]
 \leq
 \beta_{\mathsf{CPA}}(\cB_{\cA,q})+
 \frac{\sqrt{qn}}{2\gamma}.
\end{equation}
The expression \(\beta_{\mathsf{CPA}}\) is the supplied winning-probability
bound for the underlying scheme.
The formal theorem is given in the artifact as
\texttt{NoiseFloodingSecure.is\_secure}.

The result covers adaptive oracle programs rather than adversaries prewritten as \(q\) separate query phases.
Its proof isolates one lossy transition: real flooded decryption is replaced by flooding around the plaintext recorded in the challenge table.
Every other change---opening and relinking packages, compiling calls, replacing unreachable residual decryptions, and identifying the reduction game---is proved exact.

\subsection{Why a program logic?}

SSProve's original architecture connects two proof styles: high-level equational reasoning about cryptographic packages and low-level relational reasoning about their procedures~\cite{ssprove}.
Noise flooding needs a new quantitative connection of the same kind.
The local fact is a KL-divergence bound between two implementations of one decryption call; the security claim is an additive bound between two games containing an adaptive adversary.

An ordinary approximate-equivalence judgment does not preserve the required parameter dependence.
Its sequence rule adds scalar errors, so applying Pinsker at every oracle call turns a one-call cost \(\sqrt{\epsilon/2}\) into \(q\sqrt{\epsilon/2}\).
Our Pythagorean judgment instead records a tuple of conditional KL costs.
Sequencing concatenates tuples, and a checked Micciancio--Walter rule converts the accumulated tuple to additive error only once.
After \(q\) calls of KL cost \(\epsilon\), this gives \(\sqrt{q\epsilon/2}\).

The remaining obstacle is adaptivity: a concrete SSProve adversary is one stateful oracle program, not a syntactic list of \(q\) calls.
We therefore verify a compiler that exposes selected oracle calls and prove a generic local-to-adaptive lifting theorem.
Together, the judgment and compiler form a reusable program-logic interface: prove one invariant-preserving oracle rule, then obtain a whole-program additive bound for arbitrary well-typed adaptive code.
Noise flooding is the first substantial application of this interface; \Cref{thm:compile-replace} states the general theorem.

\subsection{Contributions}

We separate the cryptographic result from the machinery used to establish it.

\emph{A checked noise-flooding reduction.} We define the IND-CPA and IND-CPAD games, the noise-flooded construction, and a concrete reduction as SSProve packages.
Rocq checks the complete hybrid chain and the loss in \Cref{eq:headline}.
The paper presents the game argument in cryptographic terms; selected formal definitions appear in the appendices.

\emph{A relational program logic for SSProve.}
We define \emph{additive-error} and \emph{Pythagorean} judgments, and prove
several inference rules in Rocq.
The logic combines unary invariant reasoning,
direct coupling and additive-error rules, KL sampling rules, tuple-concatenating
sequencing, and a final bridge to statistical distance.
% Missing
% subdistribution mass is an explicit outcome, so assertions cannot make a
% comparison spuriously easier.

\emph{A verified local-to-adaptive oracle rule.} We verify a compiler that exposes selected calls in an arbitrary concrete SSProve program.
Its main theorem lifts one invariant-preserving Pythagorean oracle judgment to a whole-program additive-error judgment with loss \(\sqrt{q\norm{s}_1/2}\).
This is a program-logic result independent of the particular FHE game.

\emph{Probability and discrete-Gaussian analysis.} We verify KL divergence and its finiteness obligations, Pinsker's inequality, a conditional-coordinate preservation theorem, and the equal-variance KL formula for integer discrete Gaussians.
These results provide the semantic soundness argument and the derived vector-Gaussian sampling rule used by the logic.

\emph{An auditable theorem boundary.} The 27,986-line artifact is organized so that games, assumptions, loss, and final theorem can be reviewed separately from proof scripts.
We report its trusted base, distinguish inherited library assumptions from local proof obligations, and provide a verified replacement for the one unverified real analysis lemma inherited through the current SSProve stack.

\subsection{Theorem boundary and conventional obligations}
\label{sec:scope}

The formally verified theorem roughly says that, given any \emph{approximately correct} and \emph{IND-CPA secure} FHE scheme, noise flooding can transform such a scheme to become IND-CPAD secure.
In particular, we assume the underlying FHE scheme includes deterministic decryption, approximate correctness with certainty (i.e. no bad key pairs exist), and a particular metric space structure for plaintexts.
Our formulation cleanly isolates one step in the entire FHE security analysis.
For an imperfectly correct scheme such as CKKS, one must add a standard game hop to capture the corresponding correctness error; this does not alter our Pythagorean hop.
We have not yet formalized this additional ``up-to-bad" game hop.

For the current stage of the project, we also assume the message space is locally isometric to $\bbZ^n$.
This captures the usual instantiation (i.e. polynomial rings over ideals) while being simpler to audit.

We have not yet formally verified any concrete FHE scheme to be IND-CPA secure or approximately correct.
Instead, we discuss the CKKS scheme in \Cref{sec:limitations}, and leave its formalization and proof engineering for future work.
As a result, we do not claim a verified deployment or concrete parameter set.

Our result also follows SSProve's usual convention: probabilistic polynomial time is not formalized in the Rocq/SSProve semantics~\cite{ssprove}, and should be checked by the usual external audit.
For the LMSS noise flooding, the standard reduction runs \(\cA\) once and adds oracle forwarding, table bookkeeping, and Gaussian sampling; no unusual complexity issue is hidden here.

\subsection{Road map}

\Cref{sec:security-model} defines the construction, attack game, and theorem.
\Cref{sec:verified-reduction} presents the game hops.
\Cref{sec:pyth-toolkit} presents the program logic and its local-to-adaptive lifting theorem, and \Cref{sec:artifact} gives the audit map and trusted base.


\Ethan{Alex's intro version below}
\section{Introduction}
\label{sec:introduction}

\begin{figure*}[t]
\centering
\begin{tikzpicture}[
  x=1cm,
  y=1cm,
  >={Latex[length=2mm]},
  dependency/.style={->,semithick,draw=black!65},
  module/.style={
    draw=black!65,
    rounded corners=2pt,
    align=center,
    text width=3.0cm,
    minimum height=1.35cm,
    inner sep=3pt,
    font=\footnotesize
  },
  interface/.style={module,fill=orange!10},
  application/.style={module,fill=green!10},
  reusable/.style={module,text width=3.35cm,fill=blue!8},
  result/.style={module,fill=violet!12,draw=black!80,very thick},
  lane/.style={font=\sffamily\footnotesize\bfseries,anchor=west}
]

% Application-specific path.
\node[lane] at (0,5.45) {APPLICATION-SPECIFIC SECURITY PROOF};

\node[interface] (interfaces) at (1.55,4.35) {%
  \textbf{Scheme and security interfaces}\\[0.2ex]
  \(S\), charts, correctness, and games};

\node[application] (games) at (5.15,4.35) {%
  \textbf{Noise-flooded construction and games}};

\node[application] (local) at (8.75,4.35) {%
  \textbf{One-call Gaussian bound and table invariant}};

\node[application] (adaptive) at (12.35,4.35) {%
  \textbf{Exact operation bridges and adaptive lossy hop}};

\node[result] (final) at (15.95,4.35) {%
  \textbf{Game reduction and checked theorem}\\[0.2ex]
  \(\beta_{\mathsf{CPA}}(\mathcal B)+
    \sqrt{qn}/(2\gamma)\)};

\draw[dependency] (interfaces) -- (games);
\draw[dependency] (games) -- (local);
\draw[dependency] (local) -- (adaptive);
\draw[dependency] (adaptive) -- (final);

% Reusable mathematical and program-logic path.
\node[lane] at (0,2.55) {REUSABLE FOUNDATIONS};

\node[reusable] (kl) at (1.85,1.35) {%
  \textbf{KL foundations}\\[0.2ex]
  Pinsker and the conditional-coordinate chain bound};

\node[reusable] (gaussian) at (6.15,1.35) {%
  \textbf{Discrete-Gaussian analysis}\\[0.2ex]
  Normalization, moments, and exact KL};

\node[reusable] (logic) at (10.45,1.35) {%
  \textbf{Semantic program logics}\\[0.2ex]
  Additive-error and Pythagorean rules};

\node[reusable] (compiler) at (14.75,1.35) {%
  \textbf{Selected-call compiler}\\[0.2ex]
  Adaptive replacement\\
  theorem};

\draw[dependency] (kl) -- (gaussian);
\draw[dependency] (gaussian) -- (logic);
\draw[dependency] (logic) -- (compiler);

% The two points at which the reusable layer discharges the application.
\draw[dependency] (gaussian.north) -- (local.south);
\draw[dependency] (logic.north) -- (local.south);
\draw[dependency] (compiler.north) -- (adaptive.south);

\end{tikzpicture}
\caption{Theory map of the Rocq development.  Each box is a conceptually
significant module or a small, tightly coupled module group; an arrow points
from a dependency to code that uses it.  The upper lane specializes the
reusable lower lane to noise flooding and culminates in the exported security
theorem.  Orange marks input interfaces, green the application-specific
security proof, blue reusable foundations, and purple the exported result.
Utility and adapter modules, repeated transitive imports, and proof-engineering
dependencies are omitted.}
\label{fig:theory-map}
\end{figure*}


Security proofs for adaptive cryptographic systems often begin with a local distributional replacement: one oracle answer in the real game is replaced by a nearby simulated answer. The difficulty is preserving the correct concrete loss when this replacement is made repeatedly and adaptively. Suppose that every call has conditional Kullback--Leibler (KL) cost at most $\varepsilon$. Applying Pinsker's inequality separately at each of $q$ calls and then using the triangle inequality gives an additive loss of
\[
q\sqrt{\varepsilon/2}.
\]
Accumulating the conditional KL costs across the complete interaction and applying Pinsker only once instead gives
\[
\sqrt{q\varepsilon/2}.
\]

This distinction is parameter-critical in the noise-flooding defense for approximate homomorphic encryption. In that application, the one-query KL cost is
\[
\varepsilon_{\mathrm{nf}}=\frac{n}{2\gamma^{2}},
\]
where $n$ is the noise-coordinate dimension and $\gamma$ controls the flooding width. The two composition strategies therefore give losses proportional, respectively, to
\[
\frac{q\sqrt{n}}{2\gamma}
\qquad\text{and}\qquad
\frac{\sqrt{qn}}{2\gamma}.
\]
Because the flooding noise is added directly to the decrypted result, this is also the difference between a flooding width growing linearly and one growing as the square root of the number of queries. The sharper composition argument therefore directly preserves numerical precision.

Fully homomorphic encryption (FHE) allows a client to outsource computation on encrypted data without revealing either the plaintext or the secret key \cite{rad78,gentry09}. Modern FHE schemes support exact modular arithmetic, Boolean computation, or approximate numerical computation \cite{fv12,tfhe,ckks}. Approximate schemes, exemplified by CKKS, deliberately permit controlled numerical error in order to efficiently evaluate arithmetic over real- or complex-valued data \cite{ckks}. This makes them attractive for data-analysis and machine-learning workloads, and CKKS is implemented by many major libraries supporting approximate arithmetic \cite{heprofiler}.

Approximation, however, creates an additional information channel. In an interactive setting, a server may obtain decryptions of homomorphic computations whose mathematical results it already knows. By comparing the expected result with the approximate value returned by the client, the server can observe the residual decryption error. Li and Micciancio showed that this error can depend on the secret key strongly enough to permit key recovery \cite{lm}.

Li and Micciancio captured the missing security guarantee through the IND-CPAD game. The game extends IND-CPA with homomorphic evaluation and restricted decryption access: a ciphertext may be decrypted only when it was produced through the legitimate game operations and when its recorded plaintexts in the two challenge worlds agree. For a perfectly correct exact scheme, such a decryption reveals nothing beyond the already known plaintext. Approximate decryption additionally reveals its residual error.

Li, Micciancio, Schultz-Wu, and Sorrell (LMSS) subsequently showed that this channel can be masked by \emph{noise flooding} \cite{lmss}. Before releasing an approximate decryption, the client adds fresh Gaussian noise whose width is chosen relative to the public error bound associated with the ciphertext. The fresh noise hides the original secret-key-dependent error, but it also consumes precision. A useful defense must therefore control not only whether the resulting distributions are close, but also the exact dependence of the required noise width on the number of adaptive queries \cite{costache2023precision,bergamaschi2025revisiting}.

For one permitted decryption query, the LMSS analysis compares two answer distributions. The real oracle adds fresh noise around the scheme's approximate decryption, whereas the simulated oracle adds the same form of noise directly around the plaintext recorded in the challenge table. Approximate correctness bounds the displacement between these two centers. An equal-variance discrete-Gaussian calculation then bounds the conditional KL divergence of the two answers by $\varepsilon_{\mathrm{nf}}$.

The remaining step is to lift this one-call statement to the adversary's complete adaptive interaction. LMSS use a composition principle studied more generally by Micciancio and Walter \cite{mw17}: retain the conditional KL costs across the transcript, add them, and convert to statistical distance only once. Micciancio and Walter call the resulting root-sum-of-squares behavior \emph{Pythagorean probability preservation}. Their original motivation includes replacing ideal samplers with finite-precision implementations inside cryptographic algorithms. In noise flooding, the same principle connects the local Gaussian calculation to the complete IND-CPAD game.

At paper level, the adaptive part of the argument can be summarized in one sentence: condition on the previous transcript, apply the one-query KL bound at every call, invoke the KL chain rule, and apply Pinsker once. Mechanizing this instruction exposes a more substantial program-level problem. The adversary's queries do not form a fixed list of $q$ experiments. The adversary is one stateful oracle program whose next call, query argument, and control flow may depend on every preceding answer. Decryption calls may be interleaved with encryption, evaluation, heap operations, random sampling, and aborting branches. Moreover, the cryptographic invariant that justifies the local correctness bound must remain true after every reachable interaction.

A machine-checked proof must therefore solve three related problems. It must retain conditional KL information rather than converting each local comparison immediately to additive error. It must expose the dynamically occurring calls to the selected oracle operation inside an adaptive program. Finally, it must thread the cryptographic invariant through the exposed calls and all intervening computation.

We solve these problems using Rocq and SSProve \cite{ssprove}. First, we define a \emph{Pythagorean relational judgment} over SSProve semantics. The judgment compares two programs while carrying a common invariant and a tuple of conditional KL budgets. Its sequence rule concatenates these tuples instead of adding their square roots, and a verified Micciancio--Walter rule performs a single final conversion to additive error. Second, we verify a trace compiler that exposes the first $q$ calls to a selected oracle operation in a well-typed SSProve program. Compilation is exact when both sides use the same oracle implementation. Combining this exactness result with the Pythagorean judgment yields a generic local-to-adaptive replacement theorem with loss
\[
\sqrt{\frac{q\lVert\mathbf{s}\rVert_{1}}{2}},
\]
where $\mathbf{s}$ is the conditional-KL budget tuple for one selected call.

The KL chain rule, Pinsker's inequality, and the underlying Pythagorean distribution theorem are known mathematical results \cite{mw17}. Our contribution is their verified integration with the semantics of concrete adaptive SSProve programs, including a selected-call transformation, invariant reasoning, aborting executions, and the connection back to game-based cryptographic security. Noise flooding is the main checked cryptographic instantiation of this general machinery.

\subsection{Result and Scope}
\label{sec:intro-result-scope}
\label{sec:scope}

Let $S$ be an approximate homomorphic-encryption scheme satisfying the scheme, correctness, noise-coordinate, and IND-CPA interfaces specified in \Cref{sec:construction}. Let $\operatorname{NF}_{\gamma}[S]$ denote its noise-flooded transform, where $\gamma>0$ is the Gaussian-width multiplier, and let $n$ be the dimension of the integer noise-coordinate representation. For every well-typed IND-CPAD adversary $\mathcal{A}$ and every decryption-query budget $q$, the formal development constructs an IND-CPA adversary $\mathcal{B}_{\mathcal{A},q}$ and proves
\begin{equation}
\begin{aligned}
&\operatorname{WinProb}\!\left[
\operatorname{IND\text{-}CPAD}_{
\operatorname{NF}_{\gamma}[S],q
}(\mathcal{A})
\right]
\\
&\qquad\leq
\beta_{\mathrm{CPA}}\!\left(
\mathcal{B}_{\mathcal{A},q}
\right)
+
\frac{\sqrt{qn}}{2\gamma}.
\end{aligned}
\label{eq:intro-main-bound}
\end{equation}
Here $\beta_{\mathrm{CPA}}(\mathcal{B}_{\mathcal{A},q})$ bounds the underlying scheme's IND-CPA winning probability. Under the convention that advantage is excess winning probability over $1/2$,
\[
\beta_{\mathrm{CPA}}(\mathcal{B})
=
\frac{1}{2}
+
\operatorname{Adv}_{\mathrm{CPA}}(\mathcal{B}).
\]
With this convention, the Rocq development exports the theorem
\mbox{\texttt{NoiseFloodingSecure.is\_secure}}.

The proof isolates one statistical transition. In the real game, a permitted decryption query is answered by adding fresh noise around the scheme's approximate decryption. In the simulated game, the same form of noise is added around the common plaintext stored in the challenge table. Once answers are centered at that stored plaintext, they no longer depend on the secret key or the hidden challenge bit, and the remaining game is implemented exactly by the IND-CPA reduction. Thus, the only paid game hop is the adaptive replacement of real flooded decryption by plaintext-centered flooded decryption.

The theorem covers stateful adversaries that may adaptively interleave encryption, evaluation, and decryption calls; it does not require the adversary to be presented as $q$ separate query phases. The game itself enforces the budget $q$ on decryption-oracle calls along nonaborting executions. The trace compiler identifies the selected calls semantically within the adversary's control flow and preserves all dependencies between successive queries.

The checked result deliberately isolates the good-execution core of the LMSS argument. It assumes deterministic decryption, probability-one---or support-level---approximate correctness, and the integer-coordinate interface specified in \Cref{sec:construction}. A concrete CKKS instantiation would additionally require an up-to-bad game hop accounting for nonzero correctness error, as well as proofs that the chosen CKKS instantiation satisfies the scheme, coordinate, and IND-CPA assumptions. We do not claim a verified CKKS implementation or a concrete CKKS parameter set.

As in SSProve generally, probabilistic polynomial running time is not represented in the Rocq semantics and remains an external complexity audit \cite{ssprove}. The explicit reduction runs $\mathcal{A}$ once and adds oracle forwarding, table bookkeeping, and discrete-Gaussian sampling; it introduces no unusual complexity transformation.

\subsection{From a One-Call KL Bound to an Adaptive Program}
\label{sec:intro-local-to-adaptive}

Machine-checked game-based proofs commonly represent challengers and adversaries as probabilistic programs. CertiCrypt established this approach inside Coq, EasyCrypt developed a dedicated language with substantial automation, and SSProve provides a modular package-based framework inside Rocq \cite{certicrypt,easycrypt,ssprove}. In SSProve, a game is assembled from packages whose procedures sample randomness, access state, and call other procedures. Package-level game transformations are connected to lower-level judgments about the procedures implementing their oracles.

The noise-flooding reduction requires a quantitative connection not supplied by an ordinary additive-error judgment. Suppose a one-call comparison has KL cost $\varepsilon$. Converting that comparison immediately by Pinsker gives additive error $\sqrt{\varepsilon/2}$. An ordinary sequencing rule then sums this scalar error across $q$ calls and obtains $q\sqrt{\varepsilon/2}$. Recovering the square-root loss requires the proof system to retain the conditional KL costs through the entire adaptive interaction.

Our Pythagorean judgment provides this information. Semantically, it represents the two program executions by transcript distributions and records a nonnegative KL budget for each exposed transcript coordinate. It also records a common unary invariant satisfied by the non-failing outcomes on both sides. The sequence rule concatenates the budget tuples of two consecutive program fragments. Only after the complete transcript has been assembled does the Micciancio--Walter rule convert the accumulated budget to an additive-error judgment. For $q$ calls, each carrying the tuple $\mathbf{s}$, this produces the loss
\[
\sqrt{\frac{q\lVert\mathbf{s}\rVert_{1}}{2}}.
\]

The judgment is defined over completed SSProve semantics so that assertion failure remains observable rather than silently disappearing as missing subdistribution mass. We prove the soundness of its sampling, reflexivity, sequencing, shared-prefix, and final-conversion rules in Rocq. Unary Hoare reasoning handles deterministic computation and invariant preservation, while only the probabilistic fragment whose distribution changes consumes KL budget.

The second obstacle is that the adversary's selected calls are semantic events rather than a fixed syntactic sequence. An SSProve adversary is one stateful abstract syntax tree. After receiving an oracle answer, it may update its state, branch, call a different oracle, sample randomness, or terminate. A proof cannot simply index into a pre-existing list of decryption phases.

We therefore verify a compiler that exposes the first $q$ calls made by a program to a chosen oracle operation. Informally, one would run the adversary until its next selected call and suspend the remainder of the computation. Directly storing such a suspension would require serializing arbitrary Rocq continuations. The compiler instead records an execution trace containing the information needed to replay the intervening computation and reconstruct the residual program. Its same-package correctness theorem proves that exposing calls and answering them with the original implementation does not change the linked program's behavior.

Combining compiler exactness with the Pythagorean sequence rule yields the generic adaptive replacement theorem. Its interface is simple: establish an invariant-preserving Pythagorean judgment for one selected oracle call, and prove that the residual package preserves the invariant for the remaining operations. The compiler and logic then produce a whole-program additive bound for every well-typed program satisfying the theorem's validity, footprint, and memory-separation conditions.

In the noise-flooding application, the invariant records initialized and related keys, agreement of the public data and challenge bit, validity of the challenge table under encryption and evaluation, and agreement of the query counters. In particular, every reachable table row carries an error bound sufficient to compare the recorded plaintext with the approximate decryption of its ciphertext. A selected decryption body consists of a shared deterministic prefix, one Gaussian-sampling fragment, and shared deterministic postprocessing. These three parts contribute the one-call tuple
\[
\left(0,\varepsilon_{\mathrm{nf}},0\right),
\qquad
\varepsilon_{\mathrm{nf}}
=
\frac{n}{2\gamma^{2}}.
\]

The compiler exposes the successive decryption calls, the invariant makes the same local Gaussian argument available after every reachable history, and the Micciancio--Walter rule is applied once after the accumulated KL budget reaches $q\varepsilon_{\mathrm{nf}}$. This yields precisely the loss in Equation~\eqref{eq:intro-main-bound}.

The resulting proof architecture separates three concerns that are intertwined in the informal reduction: the probability calculation for one oracle answer, the adaptive control flow of the adversary, and the cryptographic invariant that makes the local calculation applicable. The general theorem is not specific to FHE. It can be reused whenever an adaptive cryptographic program invokes an implementation that is conditionally KL-close to an ideal oracle, provided that an appropriate one-call rule and invariant can be established.

\subsection{Contributions}
\label{sec:intro-contributions}

We separate the general program-logic contribution from its cryptographic instantiation.

\paragraph{A Pythagorean relational judgment for SSProve.}
We define a semantic judgment that compares two probabilistic programs while retaining a tuple of conditional KL budgets and a common invariant. We verify its principal inference rules, including sequencing by tuple concatenation and a final Micciancio--Walter conversion to additive error. The underlying probability development includes finite-KL infrastructure, Pinsker's inequality, conditional-coordinate KL composition, and the treatment of zero-mass prefixes and aborting SSProve executions.

\paragraph{A verified selected-call compiler and local-to-adaptive theorem.}
We construct a trace compiler that exposes the first $q$ calls to a selected oracle operation within a concrete stateful SSProve program. Its same-package theorem establishes exactness when the exposed and residual calls use the original implementation. Combining the compiler with a one-call Pythagorean judgment yields a generic whole-program replacement theorem with loss
\[
\sqrt{\frac{q\lVert\mathbf{s}\rVert_{1}}{2}}.
\]
This theorem is generic in the program, selected operation, oracle packages, invariant, and one-call budget tuple.

\paragraph{A checked noise-flooding reduction.}
We formalize the IND-CPA and IND-CPAD experiments, the noise-flooded transform, and the concrete IND-CPA reduction as SSProve packages. We verify the invariant connecting reachable challenge-table entries to the scheme's correctness guarantee, prove the one-call discrete-Gaussian KL bound, instantiate the generic adaptive theorem with $\left(0,\varepsilon_{\mathrm{nf}},0\right)$, and check the exact endpoint transformations. The final theorem establishes Equation~\eqref{eq:intro-main-bound}.

\paragraph{An explicit and auditable theorem boundary.}
We organize the artifact so that the games, scheme assumptions, probability results, compiler theorem, reduction, and final bound can be inspected separately. We distinguish the abstract scheme, correctness, coordinate, flooding, and IND-CPA inputs from the probability and program-logic results proved inside the development. We also report the trusted dependency boundary and the status of the inherited MathComp Analysis placeholder, for whose mathematical statement the artifact provides an independent proof.

\paragraph{Organization.}
\Cref{sec:construction} defines the noise-flooded construction and the abstract assumptions used by the theorem. The complete IND-CPA and IND-CPAD experiments and reduction pseudocode appear in \Cref{app:security-games}. \Cref{sec:verified-reduction} presents the cryptographic game sequence and isolates the adaptive statistical hop. \Cref{sec:pyth-toolkit,sec:trace-compiler} develop the Pythagorean judgment, its probability-theoretic foundation, the trace compiler, and the generic local-to-adaptive theorem; \Cref{sec:closing-hop} then instantiates them to close the noise-flooding reduction. \Cref{sec:artifact} describes the artifact, theorem-audit path, and trusted base. \Cref{sec:related-work} discusses related work, and \Cref{sec:limitations} concludes.


\section{Construction, Security Setting, and Result}
\label{sec:security-model}

This section defines the construction and states the result proved later,
explaining formalization-specific conventions as they arise.
All distributions are discrete.
Games return a Boolean, and
\(\WinProb[G]=\Pr[G=\mathsf{true}]\).
We write
\(\TVD(P,Q)=\tfrac12\sum_x|P(x)-Q(x)|\) for statistical distance and
\(\KL{P}{Q}\) for KL divergence when the defining series is summable and
\(P\) is absolutely continuous with respect to \(Q\).

\subsection{Approximate FHE and its flooded transform}

CKKS introduced homomorphic arithmetic with approximate plaintext
semantics~\cite{ckks}.
LMSS subsequently gave a generic approximate-correctness formulation in which
ciphertexts carry public error estimates~\cite{lmss}.
Our formal scheme interface is a typed, executable presentation of that
setting:
\[
 S=(\mathsf{KeyGen},\mathsf{Enc},\mathsf{Eval}_1,
       \mathsf{Eval}_2,\mathsf{Dec}).
\]
For bookkeeping, the formal model represents a ciphertext as either
\(\mathsf{None}\), for an unsupported evaluation, or
\(c=\mathsf{Some}(\bar c,e)\), pairing a ciphertext with its public error bound
\(e\in\bbN\).
This option wrapper belongs to the formal API, not to the cryptographic
construction.
Unary and binary evaluation gates suffice to describe arbitrary circuits while
keeping the package interface small.

% IEEEtran can place a double-column float only after it has been encountered.
% Keep this declaration early enough for the table to appear at the top of the
% next page; the FloatBarrier at the end of this section then causes no gap.
\begin{table*}[t]
\centering
\caption{Formal theorem boundary: inputs and obligations of the checked
security theorem.
The first
five rows are inputs to the Rocq theorem; the last row explains what its
semantic model does not express.}
\label{tab:theorem-boundary}
\begin{tabular}{@{}p{.18\textwidth}p{.34\textwidth}p{.39\textwidth}@{}}
\toprule
Input & What Rocq assumes & Consequence used by the proof \\
\midrule
Scheme \(S\) & Typed, distribution-valued algorithms; key generation returns
with probability one & Defines the IND-CPA game and noise-flooded scheme \\
Noise coordinates & Dimension \(n\), chart maps, and the three laws in
\Cref{eq:charts} & Transports product discrete-Gaussian noise from
\(\bbZ^n\) to messages and measures the displacement between its centers \\
Probability-one correctness & Deterministic decryption; every output having
nonzero probability satisfies \Cref{eq:approx-correctness} & Every reachable
table ciphertext is within its recorded error bound \\
IND-CPA security & For every adversary \(B\) with the required oracle interface,
\(\WinProb[\mathsf{IND\text{-}CPA}_S(B)]\leq\beta_{\mathsf{CPA}}(B)\) &
Supplies the bound for the final game \\
Flooding parameter & \(\gamma>0\) & Makes division and the Gaussian loss
well defined \\
Adversary model & A well-typed SSProve oracle program; running time is not part
of the semantics & Separate inspection that the explicit reduction preserves
probabilistic polynomial time~\cite{ssprove} \\
\bottomrule
\end{tabular}
\end{table*}

Although the interface permits randomized decryption, the theorem specializes
it to a deterministic value \(\mathsf{dec}_0(sk,c)\) and assumes
probability-one approximate correctness.
For every encryption or evaluation output
\(c=\mathsf{Some}(\bar c,e)\) having nonzero probability and representing
\(m\),
\begin{equation}
\label{eq:approx-correctness}
 d(\mathsf{dec}_0(sk,c),m)\leq e.
\end{equation}
The artifact calls this \emph{support-level} correctness; in cryptographic
terms, the bound holds with probability one.
Nonzero correctness error would require the additional up-to-bad hop discussed
in \Cref{sec:scope}.

LMSS's noise-flooding defense postprocesses approximate decryption with fresh
Gaussian noise~\cite{lmss}.
On a concrete plaintext ring, adding such noise to a plaintext is an ordinary
algebraic operation.
For an abstract message type \(M\), however, a product discrete Gaussian is
supported on \(\bbZ^n\), not on \(M\); it cannot be used as message-valued
noise without first supplying integer coordinates.

Our formalization therefore packages the metric and the coordinates needed
for noise flooding as one interface.
The message space carries a distance \(d:M\times M\to\bbN\), and for each
center \(a\in M\) it has maps
\[
 I_a:M\rightarrow\bbZ^n,
 \qquad J_a:\bbZ^n\rightarrow M
\]
satisfying
\begin{equation}
\label{eq:charts}
 I_a(a)=0,
 \quad d(a,b)=\norm{I_a(b)}_\infty,
 \quad J_b(v)=J_a(v+I_a(b)).
\end{equation}
Here \(I_a(b)\) is the integer displacement from \(a\) to \(b\), while
\(J_a(v)\) interprets the integer displacement \(v\) as a message based at
\(a\).
Thus \(J_a\) transports a distribution on \(\bbZ^n\) to a distribution on
messages.
The three laws say that the coordinates are centered, that their
\(\ell_\infty\) norm agrees with message distance, and that changing the base
point translates the integer coordinates.
In particular, they let the proof compare flooding around two messages as
centered and shifted Gaussians in one common copy of \(\bbZ^n\).

For \(\gamma>0\), noise flooding around \(x\in M\) is consequently the
following single construction:
\begin{equation}
\label{eq:flood}
\begin{aligned}
 \mathsf{Flood}_e(x):\quad
 \sigma_e&=\max\{1,e\}\gamma,\\
 \eta&\leftarrow D_{\bbZ^n,0,\sigma_e^2},\\
 \mathsf{return}\;&J_x(\eta).
\end{aligned}
\end{equation}
Here \(D_{\bbZ^n,0,\sigma_e^2}\) is the product of \(n\) independent
integer discrete Gaussians.
The sampling step lives in \(\bbZ^n\), and \(J_x\) makes its result a message;
an arbitrary metric space without such maps would not support
\Cref{eq:flood}.
This proof-specific interface is easier to audit than committing the generic
theorem to a concrete CKKS polynomial quotient ring.
The noise-flooded transform \(\mathsf{NF}_{\gamma}[S]\) keeps key generation,
encryption, and evaluation unchanged and decrypts a valid tagged ciphertext by
\[
 a\leftarrow\mathsf{Dec}(sk,c);
 \quad \mathsf{Flood}_e(a).
\]
On input \(\mathsf{None}\), decryption fails; SSProve represents this by a
subdistribution of total weight zero.

We use the standard multi-query left-or-right IND-CPA experiment and the
IND-CPAD experiment of Li and Micciancio~\cite{lm}; their full definitions
appear in \Cref{alg:ind-cpa-game,alg:ind-cpad-game,alg:ind-cpad-oracles}.
For the discussion below, the essential difference is that IND-CPAD adds
evaluation and restricted decryption: a decryption query is answered only
when the two plaintexts recorded for that ciphertext agree.
We write the assumed IND-CPA bound as
\[
 \WinProb[\mathsf{IND\text{-}CPA}_{S}(\cB)]
 \leq \beta_{\mathsf{CPA}}(\cB).
\]

\subsection{The one-query decryption replacement}
\label{sec:simulation-target}

The reduction compares two ways to answer a decryption query when the two
recorded plaintexts agree.
Fix a row that can occur with nonzero probability,
\[
 T_i=(m,m,c),
 \qquad c=\mathsf{Some}(\bar c,e).
\]
Under deterministic decryption, the real and simulated answers are
\begin{align}
\label{eq:real-sim-decrypt}
 \mathsf{ODec}_{\mathsf{real}}(i):\quad&
   y\leftarrow\mathsf{Flood}_e(\mathsf{dec}_0(sk,c));
   \ \mathsf{return}\ \mathsf{Some}(y),\notag\\
 \mathsf{ODec}_{\mathsf{sim}}(i):\quad&
   y\leftarrow\mathsf{Flood}_e(m);
   \ \mathsf{return}\ \mathsf{Some}(y).
\end{align}
Both versions perform the same bounds checks and counter updates, refuse the
same queries, and abort in the same circumstances.
The simulated answer needs neither \(sk\) nor \(b\): the common plaintext
\(m\) is already recorded in the table.

The probability-one correctness assumption in \Cref{eq:approx-correctness}
gives
\(d(\mathsf{dec}_0(sk,c),m)\leq e\).
By the change-of-basepoint law in \Cref{eq:charts}, the simulated answer can
be expressed in the real answer's chart by shifting the integer Gaussian by
\(I_{\mathsf{dec}_0(sk,c)}(m)\).
The distance law bounds the \(\ell_\infty\) norm of this shift by \(e\), so
each of its \(n\) coordinates has magnitude at most \(e\).
The equal-variance discrete-Gaussian identity
\begin{equation}
\label{eq:dg-kl}
 \KL{\DG{\mu}{\sigma^2}}{\DG{\nu}{\sigma^2}}
   =\frac{(\nu-\mu)^2}{2\sigma^2}.
\end{equation}
After fixing the adversary's prior view, \(m\), \(c\), and \(e\) are fixed.
The identity bounds the KL divergence \(\epsilon_i\) between the two answer
distributions for this call by
\begin{equation}
\label{eq:epsilon-nf}
 \epsilon_i
 \leq
 \frac{ne^2}{2\max\{1,e\}^2\gamma^2}
 \leq
 \epsilon_{\mathsf{nf}},
 \qquad
\epsilon_{\mathsf{nf}}=\frac{n}{2\gamma^2}.
\end{equation}

This is a one-call statement conditioned on an arbitrary prior interaction
history.
The central problem is to preserve it across the adversary's \(q\) adaptive
queries without converting to statistical distance after each call.

The underlying noise-flooding reduction is due to LMSS~\cite{lmss}.
Our contribution is to express its assumptions as the explicit interface in
\Cref{tab:theorem-boundary} and check the resulting adaptive reduction in Rocq.

\begin{theorem}[Checked noise-flooding reduction]
\label{thm:main}
Let \(S\), its noise-coordinate interface, probability-one correctness proof,
IND-CPA security interface, and \(\gamma>0\) satisfy
\Cref{tab:theorem-boundary}.
For every IND-CPAD adversary \(\cA\) represented as a well-typed SSProve oracle
program, and every \(q\in\bbN\), the formalization constructs an IND-CPA adversary
\(\cB_{\cA,q}\) with the required oracle interface and proves
\begin{align*}
 &\WinProb[
   \mathsf{IND\text{-}CPAD}_{\mathsf{NF}_{\gamma}[S],q}(\cA)]\\
 &\quad\leq \beta_{\mathsf{CPA}}(\cB_{\cA,q})
       +\sqrt{q\epsilon_{\mathsf{nf}}/2}\\
 &\quad\leq \beta_{\mathsf{CPA}}(\cB_{\cA,q})
       +\frac{\sqrt{qn}}{2\gamma}.
\end{align*}
\end{theorem}

The first inequality keeps the per-query bound
\(\epsilon_{\mathsf{nf}}\) visible; substituting
\(\epsilon_{\mathsf{nf}}=n/(2\gamma^2)\) from \Cref{eq:epsilon-nf} gives the
second.
Thus IND-CPA security is an assumption on \(S\), while IND-CPAD security of its
noise-flooded version is the conclusion.

\FloatBarrier

\section{The Reduction and Its Missing Link}
\label{sec:verified-reduction}

\Cref{sec:simulation-target} isolates the local change behind the LMSS
reduction: replace \(\mathsf{ODec}_{\mathsf{real}}\) by the
plaintext-centered \(\mathsf{ODec}_{\mathsf{sim}}\).
The latter can be implemented without the secret key and hence inside an
IND-CPA reduction.
For one call, \Cref{eq:epsilon-nf} bounds the KL cost.
We now lift that local comparison to the complete adaptive experiment.

\subsection{Three games}

Fix an IND-CPAD adversary \(\cA\) making at most \(q\) decryption queries.
Consider the following games:
\begin{align*}
 G_0 &=
   \mathsf{IND\text{-}CPAD}_{\mathsf{NF}_{\gamma}[S],q}(\cA),\\
 G_1 &=
   \mathsf{IND\text{-}CPAD}^{\mathsf{sim}}_
     {\mathsf{NF}_{\gamma}[S],q}(\cA),\\
 G_2 &=
   \mathsf{IND\text{-}CPA}_{S}(\cB_{\cA,q}).
\end{align*}
Game \(G_0\) is the real experiment.
Game \(G_1\) replaces the real decryption oracle by
\(\mathsf{ODec}_{\mathsf{sim}}\) from \Cref{eq:real-sim-decrypt}.
This simulated-decryption game is the intermediate game used by
LMSS~\cite{lmss}.
Game \(G_2\) runs the concrete reduction from
\Cref{alg:ind-cpa-reduction}.

The second transition is exact.
Once decryption answers are centered at the common recorded plaintext, they no
longer require the secret key or the hidden challenge bit.
The reduction can therefore forward encryption queries to its IND-CPA oracle,
perform evaluation and table bookkeeping locally, and simulate every permitted
decryption query with the same distribution as \(G_1\).
Consequently,
\[
 \WinProb[G_1]=\WinProb[G_2].
\]

The entire cryptographic proof is thus the chain
\begin{equation}
\label{eq:conceptual-game-chain}
 G_0
 \xrightarrow{\sqrt{q\epsilon_{\mathsf{nf}}/2}}
 G_1
 \xrightarrow{0}
 G_2.
\end{equation}
The first arrow is the only paid hop.
Together with the assumed IND-CPA bound for \(G_2\), it immediately implies
\Cref{thm:main}.

\subsection{The tempting one-line proof}
\label{sec:paid-hop}

\Cref{eq:epsilon-nf} bounds the conditional KL cost of every reachable
decryption query whose two recorded plaintexts agree by
\(\epsilon_{\mathsf{nf}}\).
The two games agree when those plaintexts differ and when an assertion fails,
so those branches spend no budget.

It is tempting to finish in one sentence:
apply the KL chain rule over the \(q\) decryption interactions, obtain total KL cost
\(q\epsilon_{\mathsf{nf}}\), and apply Pinsker once.
This gives
\begin{equation}
\label{eq:central-game-hop}
 \TVD(G_0,G_1)
 \leq \sqrt{q\epsilon_{\mathsf{nf}}/2}.
\end{equation}
This is the KL-composition argument used by LMSS~\cite{lmss}.
A machine-checked proof must identify the relevant conditional distributions
and establish the one-call bound after every adversarially chosen history.

\subsection{Where adaptivity enters}
\label{sec:adaptive-invariant}

The adversary's decryption calls do not form a fixed product experiment.
After seeing an answer, \(\cA\) may change its state, choose a different oracle,
or decide whether and where to make its next decryption query.
In SSProve, \(\cA\) is one stateful program with calls embedded in its control
flow, not a list of \(q\) pre-existing query phases.
Thus ``the first \(q\) decryption calls'' is a semantic notion, and the
conditional KL bound must hold after every reachable history.

The proof also depends on a state invariant \(\Phi\) relating the real and
simulated executions.
It records initialized, good keys; agreement of their public data and hidden
challenge bits; validity of the challenge table under encryption and
evaluation; and agreement of the decryption counters.
In particular, for every reachable row
\((m_0,m_1,c)\) with \(c=\mathsf{Some}(\bar c,e)\), the plaintext on the
selected branch \(b\) satisfies
\[
 d(\mathsf{dec}_0(sk,c),m_b)\leq e.
\]
Initialization establishes \(\Phi\), while encryption, evaluation, and the
deterministic prefix of decryption preserve it.
This is what makes the one-call KL calculation available after an arbitrary
adaptive history.

A machine-checked proof of \Cref{eq:central-game-hop} therefore needs a bridge
with three properties:
\begin{enumerate}
\item it retains conditional KL costs instead of converting each call to
statistical distance;
\item it exposes successive calls of an arbitrary adaptive oracle program; and
\item it threads \(\Phi\) through the exposed calls.
\end{enumerate}
\Cref{sec:pyth-toolkit} develops the Pythagorean judgment and invariant rules.
\Cref{sec:trace-compiler} verifies the selected-call compiler and proves the
generic local-to-adaptive theorem independently of FHE.  The instantiation in
\Cref{sec:closing-hop} combines them to derive the paid arrow in
\Cref{eq:conceptual-game-chain}.

\section{A Program Logic for Adaptive Games}
\label{sec:pyth-toolkit}

SSProve combines high-level algebraic reasoning about packages with a
relational program logic for their low-level procedures~\cite{ssprove}.  Its
central bridge turns per-procedure exact judgments into perfect
indistinguishability of packages.  We retain that two-level architecture but
add a quantitative bridge needed by noise flooding: a bounded conditional KL
cost for one oracle implementation is lifted to an additive-error bound for a
complete adaptive program.  This section presents that contribution
independently of the FHE instantiation.

The development introduces no new programming language and assumes no logic
rules.  It defines semantic judgments over SSProve's existing
semantics, then verifies the soundness of every rule below.
The resulting interface has
three layers: unary Hoare reasoning establishes support invariants; an
additive-error judgment handles conventional game hops; and a Pythagorean
judgment retains conditional KL information until the end of a computation.

\subsection{SSProve syntax and semantics}
\label{sec:ssprove-substrate}

For completeness, we recall the fragment of SSProve on which our logic is
built~\cite{ssprove}.  SSProve embeds pure expressions shallowly in Rocq and
represents effects by a free monad.  Fix a result type \(A\).  Its raw code has
the following constructors:
\begin{equation}
\label{eq:ssprove-syntax}
\begin{split}
c ::= {}& \mathsf{return}(a)
 \mid \mathsf{call}(p,x,\kappa)
 \mid \mathsf{get}(\ell,\kappa)\\
 &{}\mid \mathsf{put}(\ell,v,c)
 \mid \mathsf{sample}(D,\kappa).
\end{split}
\end{equation}
Here \(a:A\).  An operation signature \(p:S\to T\) contains a procedure
identifier as well as its argument and result types; a call has \(x:S\) and
\(\kappa:T\to\mathsf{RawCode}(A)\).  A typed location
\(\ell:T_\ell\) is read into a continuation
\(\kappa:T_\ell\to\mathsf{RawCode}(A)\), or updated with \(v:T_\ell\).
Finally, sampling takes \(D\in\mathsf{SD}(X)\) and
\(\kappa:X\to\mathsf{RawCode}(A)\), where \(\mathsf{SD}(X)\) denotes discrete
subdistributions on \(X\).  Continuations are Rocq functions, so conditionals,
pure local computation, and structurally terminating loops are supplied by
the ambient language rather than by additional constructors.  Monadic bind
recursively pushes a continuation through this syntax.  We use the customary
notations
\begin{equation}
\label{eq:ssprove-notation}
\begin{aligned}
 x\leftarrow c_1;c_2
   &\equiv\mathsf{bind}(c_1,\lambda x.c_2),\\
 x\leftarrow p(a);c
   &\equiv\mathsf{call}(p,a,\lambda x.c),\\
 x\leftarrow\mathsf{get}\,\ell;c
   &\equiv\mathsf{get}(\ell,\lambda x.c),\\
 \mathsf{put}\,\ell:=v;c
   &\equiv\mathsf{put}(\ell,v,c),\\
 x\leftarrow D;c
   &\equiv\mathsf{sample}(D,\lambda x.c).
\end{aligned}
\end{equation}

An interface is a finite set of typed operation signatures.  The checked type
\(\mathsf{Code}_{L,I}(A)\) pairs raw code with a proof that every accessed
location belongs to the footprint \(L\) and every external call belongs to
the import interface \(I\).  A package \(P:I\to E\) is a finite map from the
procedure names in its export interface \(E\) to well-scoped implementations
\(S\to\mathsf{Code}_{L,I}(T)\).  Sequential composition \(P\circ Q\)
links \(P\)'s calls to procedures implemented by \(Q\).  At code level, the
essential clause is
\begin{equation}
\label{eq:ssprove-link}
\begin{aligned}
 &\mathsf{codeLink}(\mathsf{call}(p,a,\kappa),Q)\\
 &\qquad =
 x\leftarrow Q.p(a);
 \mathsf{codeLink}(\kappa(x),Q),
\end{aligned}
\end{equation}
with \(\mathsf{codeLink}\) acting recursively on the remaining constructors.
Parallel composition combines disjoint export interfaces, while the identity
package forwards every call.  Package validity tracks the import, export, and
memory-footprint side conditions used by these operations.

Each location carries a type and a default initial value.  A heap \(h\) maps
locations to values of their declared types; lookup of an unset location
returns its default.  After linking has eliminated external calls, the
denotation of closed code is a state-transforming subdistribution
\begin{equation}
\label{eq:ssprove-denotation}
 \semantic{c}:\mathsf{Heap}\longrightarrow
  \mathsf{SD}(A\times\mathsf{Heap}).
\end{equation}
Writing \(\delta_z\) for the point distribution at \(z\), its defining
equations are
\begin{align}
 \semantic{\mathsf{return}(a)}_h
   &=\delta_{(a,h)},\notag\\
 \semantic{\mathsf{get}(\ell,\kappa)}_h
   &=\semantic{\kappa(h(\ell))}_h,\notag\\
 \semantic{\mathsf{put}(\ell,v,c)}_h
   &=\semantic{c}_{h[\ell\mapsto v]},\notag\\
 \semantic{\mathsf{sample}(D,\kappa)}_h
   &=\sum_{x}D(x)\,\semantic{\kappa(x)}_h.
\label{eq:ssprove-effect-semantics}
\end{align}
The sums are pointwise sums of subdistributions.  In particular, semantic
bind threads both the returned value and the updated heap:
\begin{align}
 \semantic{x\leftarrow c;\kappa(x)}_h
 &=
 \sum_{(a,h')}
   \semantic{c}_h(a,h')\,
   \semantic{\kappa(a)}_{h'}.
\label{eq:ssprove-bind-semantics}
\end{align}
SSProve defines failure by sampling the null subdistribution, and
\(\mathsf{assert}(b)\) as \(\mathsf{return}(())\) when \(b\) holds and failure
otherwise.  Thus failed assertions have the zero subdistribution and their
continuations are not run.  A closed game is evaluated by resolving its
\(\mathsf{Run}\) procedure, starting from the default heap, and projecting
the result component of \Cref{eq:ssprove-denotation}.

\subsection{Completed semantics and three judgments}

The denotation in \Cref{eq:ssprove-denotation} may have weight below one.  For
a subdistribution \(\mu\) on \(X\), define its completion on
\(X\cup\{\bot\}\) by
\[
 \completed{\mu}(x)=\mu(x),\qquad
 \completed{\mu}(\bot)=1-\sum_x\mu(x).
\]
Completion makes assertion failure observable.
Without this, its becomes tedious (and sometimes tricky) to write down couplings for sub-distributions and then estimate the probabilities of post-conditions.

The unary judgment
\[
 \vDash\mathsf{Hoare}\{P\}\ c\ \{Q\}
\]
is the standard Hoare triple over the SSProve code.
It states that whenever the pre-condition $P$ is satisfied,
the post-condition $Q$ would be satisfied with certainty (i.e. with probability $1$).

The additive-error judgment
\[
 \vDash\AEJ\{P\}\ c_L\approx_\delta c_R\ \{Q\}
\]
states that, from any pair of initial configurations satisfying \(P\), the
completed outputs of \(c_L\) and \(c_R\) admit a coupling in which \(Q\) fails
with probability at most \(\delta\).
For equality post-conditions, this is the
ordinary statistical-distance bound.
We prove consequence, sequencing with loss \(\delta_1+\delta_2\),
triangle, projection, and direct-coupling rules.

The Pythagorean judgment
\begin{equation}
\label{eq:pyth-judgment}
 \vDash\Pyth\{P\}\ c_L\approx_s c_R\ \{Q\}
\end{equation}
captures the Micciancio-Walter \cite{mw17} Pythagorean prerequisites for KL divergence.
It uses a nonnegative tuple \(s=(s_1,\ldots,s_k)\).  For every related input pair,
it supplies two joint transcript distributions whose final
marginals are the programs' outputs.  At coordinate \(i\), for
every common prefix, the conditional KL divergence is finite and at most
\(s_i\).
Outputs on both sides additionally satisfy the post-condition \(Q\).
The transcript is semantic rather than a list of syntactic sampling sites: it may record a
whole program fragment as one coordinate, or expose multiple coordinates for
composition.

This judgment deliberately stores more information than an immediate
statistical-distance bound.  Its sequence rule concatenates tuples (i.e. $s\doubleplus s'$) rather
than adding square roots, allowing one nonlinear conversion after the full
adaptive transcript has been constructed.

\subsection{Core proof rules}
\label{sec:logic-rules}

\Cref{fig:logic-rules} shows the rules that drive the generic development.
They are schematic to suppress SSProve typing, full-mass, support, and heap
side conditions; the artifact contains their concrete statements.

\begin{figure*}[t]
\centering
\small
\begin{mathpar}
\inferrule*[right=\textsc{KL-Sample}]
  {\mathsf{finiteKL}(D_L,D_R) \\
   \KL{D_L}{D_R}\leq\epsilon \\
   \mathsf{weight}(D_L)=\mathsf{weight}(D_R)=1}
  {\vDash\Pyth\{P\}\
     x\leftarrow D_L
     \approx_{[\epsilon]}
     x\leftarrow D_R\ \{Q\}}

\inferrule*[right=\textsc{Pyth-Refl}]
  {P\Longrightarrow x_L=x_R\wedge m_L=m_R \\
   0\leq s_i\ \text{for every }i}
  {\vDash\Pyth\{P\}\ c\approx_s c\ \{Q\}}

\inferrule*[right=\textsc{Pyth-Seq}]
  {\vDash\Pyth\{P\}\ c_L\approx_s c_R\ \{M\} \\
   \vDash\Pyth\{M^{=}\}\ k_L\approx_t k_R\ \{Q\}}
  {\vDash\Pyth\{P\}\
     (x\leftarrow c_L;k_L(x))
     \approx_{s\doubleplus t}
     (x\leftarrow c_R;k_R(x))\ \{Q\}}

\inferrule*[right=\textsc{Shared-Prefix}]
  {\vDash\mathsf{Hoare}\{P_0\}\ c\ \{M\} \\
   \vDash\Pyth\{M^{=}\}\ k_L\approx_s k_R\ \{Q\}}
  {\vDash\Pyth\{P\}\
     (x\leftarrow c;k_L(x))
     \approx_{[0]\doubleplus s}
     (x\leftarrow c;k_R(x))\ \{Q\}}

\inferrule*[right=\textsc{Micciancio--Walter}]
  {\vDash\Pyth\{P\}\ c_L\approx_s c_R\ \{Q\}}
  {\vDash\AEJ\{P\}\ c_L
     \approx_{\sqrt{\norm{s}_1/2}}c_R\ \{o_L=o_R\}}
\end{mathpar}
\caption{Representative checked rules.  Here \(M^{=}\) requires equal
intermediate values and heaps satisfying \(M\), and \(\doubleplus\) denotes
tuple concatenation.  The Rocq theorems additionally expose all typing,
support, invariant, and nonnegativity premises.}
\label{fig:logic-rules}
\end{figure*}

\textsc{KL-Sample} embeds a distribution-level KL bound as a program judgment.
\textsc{Pyth-Refl} relates a program to itself at zero cost.
\textsc{Pyth-Seq} is the essential rule: its proof
constructs a joint transcript for semantic bind and proves conditional bounds
for both valid and zero-mass prefixes.  The development also contains a
shared-continuation rule, consequence rules, closed-program variants,
and the additive-error rules mentioned above.

The mixed Hoare/Pythagorean rules have been useful in practice.  Deterministic
prefixes, table operations, other oracle calls, and postprocessing are proved to
preserve an invariant with the unary Hoare logic and inserted as zero-cost
coordinates in the Pythagorean judgment.
Only the sampling fragment whose distribution changes consumes
KL budget.  In \Cref{lem:one-call}, this decomposition produces the tuple
\((0,\epsilon_{\mathsf{nf}},0)\).

\subsection{The Micciancio-Walter Probability Lemma}

We discuss the probabilistic Pythagorean preservation lemma used to justify the Micciancio-Walter rule in
\Cref{fig:logic-rules}.
Let \(\mathcal P,\mathcal Q\) be joint distributions on \(\Omega^k\), and
write \(\mathcal P_i\mid a\) for coordinate \(i\) conditioned on prefix
\(a\in\Omega^{i-1}\). 

\begin{theorem}[Conditional-coordinate preservation]
\label{thm:pyth-probability}
Suppose \(\mathcal P,\mathcal Q\) have weight one and, for every coordinate
\(i\) and prefix \(a\), the conditional KL divergence is finite and
\[
 \KL{\mathcal P_i\mid a}{\mathcal Q_i\mid a}\leq s_i,
 \qquad s_i\geq0.
\]
Then their final-coordinate marginals satisfy
\[
 \TVD(\mathcal{P}_k,
      \mathcal{Q}_k)
 \leq \sqrt{\frac{\sum_{i=1}^{k}s_i}{2}}.
\]
\end{theorem}

At paper level, the proof expands joint mass as prefix mass times conditional
mass and groups the log likelihood ratio by coordinate.  The resulting chain
bound gives \(\KL{\mathcal P}{\mathcal Q}\leq\sum_i s_i\).  Pinsker bounds the
transcript distance, and data processing gives the displayed final-marginal
bound.  The Rocq proof makes explicit absolute continuity, conditional
zero-mass prefixes, summability of positive and negative parts, and exchanges
between finite coordinate sums and countable support sums.

Our \textsc{Micciancio--Walter} rule instantiates this theorem with the
transcripts from \Cref{eq:pyth-judgment}, then uses maximal coupling to
obtain the additive-error equality judgment.
As with all of our inference rules,
we verify it to be sound with respect to SSProve's program semantics.

\subsection{A derived discrete-Gaussian rule}

For center \(\mu\in\bbZ\) and \(\sigma>0\), we denote the discrete Gaussian distribution by
\[
 D_{\bbZ,\mu,\sigma^2}(x)
 =\frac{\exp(-(x-\mu)^2/(2\sigma^2))}
        {\sum_{z\in\bbZ}\exp(-(z-\mu)^2/(2\sigma^2))}.
\]
The development proves normalization, translation invariance of the
normalizer, summability, the centered first moment, full support, and the
equal-variance KL identity in \Cref{eq:dg-kl}.  The calculation expands the
logarithm of the ratio; equal variances cancel both normalizers, and symmetry
of the centered distribution cancels the linear moment.

These facts yield a reusable program rule for vector sampling.  If
\(u,v\in\bbZ^n\), \(\sigma>0\), and every coordinate satisfies
\[
 \frac{(v_i-u_i)^2}{2\sigma^2}\leq\epsilon,
\]
then the $n$-fold product sampler satisfy a Pythagorean
judgment with budget \([n\epsilon]\).
% The theorem includes injectivity of the SSProve encoding and heap/postcondition premises.
The analysis first constructs
an \(n\)-coordinate transcript, proves the scalar bound coordinatewise, and
then transports and aggregates it into the program-level sampling rule.  The
noise-flooding proof instantiates \(\epsilon=1/(2\gamma^2)\), obtaining
\(\epsilon_{\mathsf{nf}}=n/(2\gamma^2)\) without introducing any unverified assumptions.

\subsection{From one oracle call to an adaptive program}
\label{sec:trace-compiler}

The rules above compose a sequence of sampling operations. An oracle program
does not provide such a sequence: the location and arguments of its next call
may depend on previous answers.  Our compiler is the verified adapter
between these views.  

Intuitively speaking, the pen-and-paper analysis would run the adversary to its next oracle call and return
its query together with a suspended execution.
In SSProve, however, a suspended program is a \emph{continuation} inside the program monad: storing it directly would
require serializing arbitrary Rocq functions into SSProve-compatible data types.
Instead, as an simpler alternative, we choose to record execution traces.
In particular, we record the results of other oracle calls, heap reads
and writes, and samples; replaying them recovers the continuation.
The same-package theorem below proves this representation exact.

Write
\(\mathsf{Comp}_q(A;P_s,P_r)\) for the program obtained by exposing the first
\(q\) calls made by \(A\) to a selected operation, resolving those calls with
\(P_s\), and linking the residual program with \(P_r\).  The compiler first
satisfies an exactness theorem.

\begin{theorem}[Same-package compilation]
\label{thm:compile-same}
Subject to package typing and memory-separation conditions,
\[
 \mathsf{Comp}_q(A;P,P)
 \quad\text{and}\quad
 \mathsf{Link}(A,P)
\]
have identical SSProve output/heap semantics.
\end{theorem}

The following theorem is the main program-logic bridge.

\begin{theorem}[Replace $q$ calls]
\label{thm:compile-replace}
Suppose two implementations \(P,P'\) of a selected operation satisfy a
one-call Pythagorean judgment with nonnegative tuple \(s\), preserve invariant
\(I\), and meet the compiler's typing, footprint, and memory-separation
conditions.  For every concrete SSProve program \(A\), replacing its first
\(q\) selected calls yields
\[
\begin{aligned}
 &\vDash\AEJ\{P_I\}\ \mathsf{Comp}_q(A;P,P)\\
 &\quad\approx_{\sqrt{q\norm{s}_1/2}}
   \mathsf{Comp}_q(A;P',P)\ \{o_L=o_R\},
\end{aligned}
\]
where \(P_I\) synchronizes inputs and heaps satisfying \(I\).
\end{theorem}
\begin{proof}[Proof sketch]
We first show that $\mathsf{Comp}_q(A;P,P)$ and $\mathsf{Comp}_q(A;P',P)$
satisfy the expected Pythagorean judgment (i.e. differ by the tuple $s^q$)
by induction on $q$ and applying our Pythagorean sequence rule.
The $q$-th call contributes \(s\) and preserves the invariant \(I\).

After \(q\) calls, the accumulated tuple has norm
\(q\norm{s}_1\); \textsc{Micciancio--Walter} is applied once, giving the
desired additive error bound.
\end{proof}

The theorem is generic in the program, packages, selected oracle,
invariant, and budget tuple.  The noise-flooding game supplies only the
one-call premise from \Cref{lem:one-call}; the program logic and compiler then
handle adversarial control flow.  Combined with
\Cref{thm:compile-same}, this is our quantitative counterpart of SSProve's
local relational reasoning for package indistinguishability.

\subsection{Why the extra judgment is necessary}

The additive-error layer remains essential for exact hops, up-to-bad
arguments, and the final triangle calculation.  Its ordinary sequence rule,
however, maps \(\delta_1,\delta_2\) to \(\delta_1+\delta_2\).  Applying Pinsker
inside each decrypt proof would therefore pay
\(q\sqrt{\epsilon_{\mathsf{nf}}/2}\).  The Pythagorean layer delays that
nonlinear conversion and obtains \(\sqrt{q\epsilon_{\mathsf{nf}}/2}\).  It is
not a new security definition; it is a verified relational abstraction that
preserves a parameter-critical cryptographic argument across adaptive,
stateful program composition.

\section{Mechanization and Trusted Base}
\label{sec:artifact}

The development contains 27,987 lines of Rocq, excluding the manuscript's
copied source excerpts.  Since a majority of the codebase was written with AI,
the line count measures the scope of the checked development rather than human
effort.
\Cref{fig:theory-map} separates the application-specific security argument
from the reusable probability and program-logic layers.

\begin{figure*}[t]
\centering
\begin{tikzpicture}[
  x=1cm,
  y=1cm,
  >={Latex[length=2mm]},
  dependency/.style={->,semithick,draw=black!65},
  module/.style={
    draw=black!65,
    rounded corners=2pt,
    align=center,
    text width=3.0cm,
    minimum height=1.35cm,
    inner sep=3pt,
    font=\footnotesize
  },
  interface/.style={module,fill=orange!10},
  application/.style={module,fill=green!10},
  reusable/.style={module,text width=3.35cm,fill=blue!8},
  result/.style={module,fill=violet!12,draw=black!80,very thick},
  lane/.style={font=\sffamily\footnotesize\bfseries,anchor=west}
]

% Application-specific path.
\node[lane] at (0,5.45) {APPLICATION-SPECIFIC SECURITY PROOF};

\node[interface] (interfaces) at (1.55,4.35) {%
  \textbf{Scheme and security interfaces}\\[0.2ex]
  \(S\), charts, correctness, and games};

\node[application] (games) at (5.15,4.35) {%
  \textbf{Noise-flooded construction and games}};

\node[application] (local) at (8.75,4.35) {%
  \textbf{One-call Gaussian bound and table invariant}};

\node[application] (adaptive) at (12.35,4.35) {%
  \textbf{Exact operation bridges and adaptive lossy hop}};

\node[result] (final) at (15.95,4.35) {%
  \textbf{Game reduction and checked theorem}\\[0.2ex]
  \(\beta_{\mathsf{CPA}}(\mathcal B)+
    \sqrt{qn}/(2\gamma)\)};

\draw[dependency] (interfaces) -- (games);
\draw[dependency] (games) -- (local);
\draw[dependency] (local) -- (adaptive);
\draw[dependency] (adaptive) -- (final);

% Reusable mathematical and program-logic path.
\node[lane] at (0,2.55) {REUSABLE FOUNDATIONS};

\node[reusable] (kl) at (1.85,1.35) {%
  \textbf{KL foundations}\\[0.2ex]
  Pinsker and the conditional-coordinate chain bound};

\node[reusable] (gaussian) at (6.15,1.35) {%
  \textbf{Discrete-Gaussian analysis}\\[0.2ex]
  Normalization, moments, and exact KL};

\node[reusable] (logic) at (10.45,1.35) {%
  \textbf{Semantic program logics}\\[0.2ex]
  Additive-error and Pythagorean rules};

\node[reusable] (compiler) at (14.75,1.35) {%
  \textbf{Selected-call compiler}\\[0.2ex]
  Adaptive replacement\\
  theorem};

\draw[dependency] (kl) -- (gaussian);
\draw[dependency] (gaussian) -- (logic);
\draw[dependency] (logic) -- (compiler);

% The two points at which the reusable layer discharges the application.
\draw[dependency] (gaussian.north) -- (local.south);
\draw[dependency] (logic.north) -- (local.south);
\draw[dependency] (compiler.north) -- (adaptive.south);

\end{tikzpicture}
\caption{Theory map of the Rocq development.  Each box is a conceptually
significant module or a small, tightly coupled module group; an arrow points
from a dependency to code that uses it.  The upper lane specializes the
reusable lower lane to noise flooding and culminates in the exported security
theorem.  Orange marks input interfaces, green the application-specific
security proof, blue reusable foundations, and purple the exported result.
Utility and adapter modules, repeated transitive imports, and proof-engineering
dependencies are omitted.}
\label{fig:theory-map}
\end{figure*}


The mechanization contains a few layers.  The scheme layer
defines encryption games and the noise-flooded transform.  The
probability layer develops facts about countable discrete distributions,
without reference to programs.  The logic layer lifts those facts to SSProve
code and proves the selected-call compiler sound.  Finally, the security layer
instantiates the generic logic with the decryption invariant, constructs the
IND-CPA reduction, and composes the exact and quantitative game transitions.
The top-level result is a functor over precisely the scheme, chart,
probability-one correctness, IND-CPA security, and flooding-parameter
interfaces described in \Cref{sec:construction}; probability lemmas and
program rules are proved dependencies.

\subsection{Formalization of judgments and rules}

We follow SSProve's convention and engineering framework: pure computation is
embedded in Rocq, while probabilistic and stateful effects are represented by
SSProve's free-monadic program syntax, packages, heaps, and linking operations
\cite{ssprove}.  The new logic does not extend Rocq with a primitive judgment
or a collection of trusted proof rules.
Judgments are formalized as ordinary Rocq propositions over the denotation \texttt{Pr\_code},
and inference rules are lemmas which prove such judgments.

For example, the exact definition of the additive-error judgment is shown
below.  Its types range over SSProve programs and heaps; \texttt{complete}
adds an explicit outcome for missing subdistribution mass.

\begingroup
\setmonofont{FreeMono}[Scale=MatchLowercase]
\inputminted[
  fontsize=\scriptsize,
  breaklines,
  linenos,
  firstnumber=57,
  firstline=57,
  lastline=69
]{coq}{rocq-excerpts/theories/ProgramLogics/Ae.v}
\endgroup

Thus, for every related pair of inputs and initial heaps, the proposition
requires a coupling of the completed output-and-heap distributions in which
the postcondition holds with probability at least \(1-\varepsilon\).
The Pythagorean judgment is encoded in the same style.  For each related input
pair it existentially quantifies two full transcript distributions, requires
their final marginals to equal the two completed program denotations, and
requires their supports to satisfy the postcondition.  Its predicate
\texttt{pythDist} additionally imposes full mass and the coordinatewise
conditional-KL bounds from \Cref{def:pyth-judgment}.
To place every result type in one transcript space, the Rocq definition
injectively encodes typed output/heap pairs in \texttt{nat * heap} and
represents \(\bot\) by \texttt{None}.

The rules in \Cref{fig:logic-rules} are proved implications between these
definitions.  For instance, the source-level statement of the final
KL-to-additive-error bridge is:

\begingroup
\setmonofont{FreeMono}[Scale=MatchLowercase]
\inputminted[
  fontsize=\scriptsize,
  breaklines,
  linenos,
  firstnumber=529,
  firstline=529,
  lastline=542
]{coq}{rocq-excerpts/theories/ProgramLogics/Pyth.v}
\endgroup

Here \texttt{pythagorean\_tv\_bound s} is
\(\sqrt{\sum_i s_i/2}\).  Its proof applies the verified probability theorem
to the transcript witnesses, projects to their final marginals, and invokes
the completed-distribution coupling theorem.  The sampling, sequencing,
mixed Hoare/Pythagorean, and compiler rules have the same status: each is a
Rocq lemma proved from the program semantics.  Consequently, a Rocq kernel check of
the security theorem also checks the soundness chain from the distribution
lemmas through the program judgments.

\subsection{Probability formalization}

The probability development works over MathComp's countable discrete
distributions, rather than restricting the Gaussian argument to finite
support.  It defines KL finiteness as absolute continuity together with
summability of the KL integrand, proves Pinsker's inequality, and proves the
conditional-coordinate chain bound used in
\Cref{thm:pyth-probability}.  It also constructs completion and maximal
couplings, so that a distribution-level total-variation bound yields the
program-level proposition above.

The distributions used here, notably the discrete Gaussians, have countably
infinite support.  Consequently, the KL chain rule involves infinite series.
A paper proof may leave their convergence implicit; the formalization must
establish it before decomposing transcript KL into conditional KL costs.  It
also treats histories of probability zero explicitly.  These results yield the
chain bound used in the square-root composition argument.

For the one-call bound, the development also proves that the discrete Gaussian
is normalized and has full support, that its normalizer is invariant under
integer shifts, and that its centered first moment is zero.  These facts yield
the exact equal-variance KL identity in \Cref{eq:dg-kl} and discharge the
Gaussian sampling rule.  Further theorem statements and proof structure
appear in \Cref{sec:verified-kl-analysis}.

\subsection{The checked theorem}

At the application boundary, the functor defines the concrete reduction
\(\cB_{\cA,q}\) and makes the assumed IND-CPA bound and the statistical loss
explicit:

\begingroup
\setmonofont{FreeMono}[Scale=MatchLowercase]
\inputminted[
  fontsize=\scriptsize,
  breaklines,
  linenos,
  firstnumber=530,
  firstline=530,
  lastline=533
]{coq}{rocq-excerpts/theories/Security/NoiseFloodingSecurity/Final.v}
\endgroup

After proving that the constructed reduction has the IND-CPA adversary
interface and deriving the completed-output inequality, the development
exports the following theorem.

\begingroup
\setmonofont{FreeMono}[Scale=MatchLowercase]
\inputminted[
  fontsize=\scriptsize,
  breaklines,
  linenos,
  firstnumber=1953,
  firstline=1953,
  lastline=1956
]{coq}{rocq-excerpts/theories/Security/NoiseFloodingSecurity/Final.v}
\endgroup

The premise says that \(A\) is a well-typed package with the required IND-CPAD
adversary interface.  Unfolding \texttt{security\_bound} and
\texttt{security\_loss} gives exactly \Cref{thm:main}; in particular, the
right-hand term is the assumed IND-CPA winning-probability bound for the
constructed reduction plus \(\sqrt{qn}/(2\gamma)\).  This source statement also
makes clear where the theorem boundary lies: scheme structure, charts,
probability-one correctness, IND-CPA security, and positivity of \(\gamma\)
enter through the enclosing functor, whereas reduction validity, the
discrete-Gaussian calculation, compiler correctness, and logic soundness have
already been proved.

\subsection{Reproducibility}

The complete artifact checks with Rocq 9.0.1, \texttt{rocq-ssprove}
development version at commit \texttt{c6d7d4bc3a}, MathComp Analysis 1.16.0,
and MathComp algebra tactics 1.2.7.  In a fresh switch on the authors' machine,
a clean build took about six minutes with four parallel jobs; this is a single
observation for reviewer planning, not a performance claim.  The manuscript's
Rocq listings are refreshed from the checked sources during its build, so the
displayed propositions and theorem statements do not silently diverge from
the artifact.

\subsection{Trusted computing base}
\label{sec:tcb}

The trusted base includes the Rocq kernel and standard library; MathComp,
MathComp Analysis, and their real-number/classical infrastructure; and
SSProve's program, heap, package, linking, and subdistribution semantics.
The scheme, chart, correctness, parameter, and IND-CPA interfaces from
\Cref{sec:construction} are theorem parameters and must be discharged by a
concrete instantiation.

The new KL, Gaussian, relational, compiler, and game-hop results contain no
local \texttt{Admitted}.  The installed dependency graph nevertheless reports
MathComp Analysis's
\texttt{realsum.\_\_admitted\_\_interchange\_psum}: current SSProve
distribution semantics are built on these MathComp definitions and reach the
upstream placeholder~\cite{ssprove}.  This is an inherited assumption, not an
unproven axiom introduced by our security development.

We independently discharge the mathematical obligation.  The artifact proves
the same statement as \texttt{interchange\_psum\_proved};
\texttt{Print Assumptions} for this replacement reports only the usual
extensionality and choice principles and not the admitted upstream lemma.  The
proof has since been merged to MathComp Analysis's development branch,
together with proofs of \texttt{psumB} and \texttt{interchange\_sup}, for
inclusion in the planned 1.17.0 release.  Until an upstream release replaces
the existing call sites, \texttt{Print Assumptions} for the top-level theorem
will continue to display the placeholder's name.

Besides library and classical assumptions, \texttt{Print Assumptions} for an
instantiated theorem reports the concrete scheme, chart, correctness, and
IND-CPA interfaces, together with positivity of the discrete-Gaussian width
multiplier \(\gamma\).  It does not report any program-logic rule, compiler
correctness theorem, or discrete-Gaussian KL identity: these are verified
inside the development.

Selected definitions of the scheme interface, games, transform,
and reduction remain in \Cref{app:formal-interface-excerpts}.  Together with
the source-level theorem above, they expose the small specification surface on
which a cryptographic audit should concentrate; the Rocq kernel checks that the
remaining proof terms connect that surface by the argument in
\Cref{sec:verified-reduction}.  As in SSProve and EasyCrypt, running time is not
part of the program semantics and must be inspected separately
\cite{easycrypt,ssprove}.  Here the
reduction runs the adversary once, forwards its oracle calls, and adds table
operations and discrete-Gaussian sampling.

\section{Related Work}
\label{sec:related-work}

\emph{Approximate-FHE attacks and defenses.}
Li and Micciancio identified the secret-dependent approximation-error channel
in CKKS-style decryption~\cite{lm}.  LMSS introduced the IND-CPAD formulation
and used differential-privacy and noise-flooding techniques to recover
security under explicit parameter conditions~\cite{lmss}.  Subsequent work
has sharpened the practical precision tradeoff~\cite{costache2023precision}
and revisited concrete noise-flooding parameters~\cite{bergamaschi2025revisiting}.
Attacks against exact FHE at aggressive parameters show that decryption-query
security is not confined to approximate schemes~\cite{cheon2024attacks}.
Our contribution is neither a new attack nor a better concrete parameter set:
it machine-checks the conditional game reduction and the square-root
composition mechanism on which such parameters rely.

\emph{Mechanized game-based cryptography.}
SSProve combines state-separating package algebra with a probabilistic
relational program logic in Rocq and validates the framework on encryption,
KEM--DEM, and sigma-protocol case studies~\cite{ssprove}.  We retain its
package semantics and security-proof style, but add quantitative judgments for
completed countable distributions and a selected-call compiler.  Recent
Nominal-SSProve work makes package state separation intrinsic and reports
errors found in an ElGamal mechanization~\cite{nominal-ssprove}; a further
development mechanizes query-counting, multi-instance, and nested hybrid
arguments~\cite{nested-hybrids}.  Those results address exact package
composition and general hybrid structure.  Our compiler instead preserves an
adaptive sequence of conditional KL costs and derives a sublinear
statistical-distance loss.  Integrating the two approaches is promising but
not part of this artifact.

EasyCrypt offers mature adversarial module reasoning; reflection techniques
can move distribution-level steps into executable proofs~\cite{reflection}.
Our abstraction boundary differs because SSProve adversaries are concrete
program syntax: the verified trace compiler factors selected calls before the
replacement theorem is applied.  We do not claim that trace compilation is
preferable to opaque adversary rules in general; it is the bridge required by
the host framework used here.

\emph{Approximate relational logics.}
Approximate relational Hoare logics and approximate liftings were developed
largely for differential privacy~\cite{aprhl,aprhl-couplings}.  They support
rich \((\varepsilon,\delta)\) reasoning, but their standard sequential rules
add privacy or additive-error parameters.  One can instantiate general
\(f\)-divergence machinery with KL divergence, yet applying Pinsker after each
program step still loses linearly.  The contribution here is the checked
combination of (i) a transcript judgment whose sequence rule concatenates
conditional KL budgets, (ii) a final-marginal bridge that converts once, and
(iii) an oracle compiler that makes that rule applicable to adaptive SSProve
games.  The resulting theorem is specialized, deliberately, to the
game-based statistical-distance conclusion used by noise flooding.

\section{Conclusion and Future Works}
\label{sec:limitations}

Our formalization deliberately factors correctness failures from the quantitative noise-flooding argument.
Its support-level hypotheses prove the good-execution case.
For an imperfectly correct FHE such as CKKS, one more standard up-to-bad game hop must be formalized before instantiation, and the probability of generating ``bad" key pairs must be carefully accounted for.

Our formalization has also left the set of plaintext abstract as a metric space with a $\bbZ^n$ chart.
For the current result, we believe this is a simpler interface to audit than the ``usual" lattice setting involving polynomial rings modulo ideals.

Modulo the above proof engineering items, we have formally verified the central LMSS game hop \cite{lmss}.
That is, an IND-CPA and approximately correct FHE scheme can be made IND-CPAD through noise flooding.
The natural next step then, is to verify that the CKKS scheme is indeed both IND-CPA and approximately correct.
That will certainly come with its own set of other interesting challenges.
For instance, such an effort will no doubt contend with the use of randomized versus deterministic rounding, as well as the (difficult to formalize) heuristics involved in the deterministic case as found ``in the wilds".

\section*{Ethics Considerations}

This work studies and verifies a published defense against an already public decryption attack.
It performs no experiments on deployed systems, users, or private data and introduces no new exploit.
Its principal risk is overconfidence: treating the conditional theorem as a verified concrete CKKS deployment could encourage unsafe parameters.
We mitigate that risk by stating the good key-pair boundary, chart and instantiation obligations, semantic model, and trusted base explicitly.


\label{page:main-matter-end}

\section*{Acknowledgment of AI Assistance}
OpenAI Codex was used to draft prose throughout this manuscript.
The Rocq development was also substantially AI-assisted.
Both have been audited and revised by the authors accordingly.

\bibliographystyle{IEEEtran} \bibliography{reference}

\appendices
\section{IND-CPA with Decryption Game Details}
\label{app:ind-cpad-game}

This appendix records the game syntax suppressed in
\Cref{sec:games}.
It is not meant to be a second proof outline; the main text reasons about the
named Coq games and their hybrid relations.

\begin{algorithm}
\caption{IND-CPAD Challenger Skeleton}
\label{alg:ind-cpad-game}
Ingredients:
\begin{itemize}
  \item An approximate homomorphic encryption scheme
    $S=(KeyGen,Enc,Eval,Dec)$.
  \item A decrypt-query bound $q\in\bbN$.
\end{itemize}
Global challenger state:
\begin{itemize}
  \item keys $pk,evk,sk$;
  \item hidden challenge bit $b\in\{0,1\}$;
  \item table $T\in(M\times M\times C)^*$ of challenge plaintext pairs and
    produced ciphertexts;
  \item decrypt counter $d$.
\end{itemize}
\begin{algorithmic}
  \Procedure{Game}{$\cA$}
    \State $(pk,evk,sk) \gets KeyGen()$
    \State $b \gets \{0,1\}$
    \State $T \gets [\,]$
    \State $d \gets 0$
    \State $b^* \gets \cA^{OEnc,OEval,ODec}(pk,evk)$
    \State \Return $(b^*=b)$
  \EndProcedure
\end{algorithmic}
\end{algorithm}

\begin{algorithm}
\caption{IND-CPAD Oracles}
\label{alg:ind-cpad-oracles}
The oracles share the challenger state from \Cref{alg:ind-cpad-game}.
\begin{algorithmic}
  \Procedure{OEnc}{$m_0,m_1$}
    \State $c \gets Enc(pk,m_b)$
    \State $T \gets T \doubleplus [(m_0,m_1,c)]$
    \State \Return $c$
  \EndProcedure
  \State
  \Procedure{OEval}{$f,\mathcal{I}$}
    \State $\vec r \gets (T_i)_{i\in\mathcal{I}}$
    \State $m_0' \gets f((r.m_0)_{r\in\vec r})$
    \State $m_1' \gets f((r.m_1)_{r\in\vec r})$
    \State $c' \gets Eval(f,(r.c)_{r\in\vec r})$
    \State $T \gets T \doubleplus [(m_0',m_1',c')]$
    \State \Return $c'$
  \EndProcedure
  \State
  \Procedure{ODec}{$i$}
    \State \textbf{assert} $d<q$
    \State $d \gets d+1$
    \State \textbf{assert} $i<|T|$
    \If{$T_i.m_0\ne T_i.m_1$}
      \State \Return $\mathsf{None}$
    \Else
      \State \Return $\mathsf{SampleCurrent}(T_i)$
    \EndIf
  \EndProcedure
\end{algorithmic}
\end{algorithm}

Here \(\mathsf{SampleCurrent}\) is instantiated by the surrounding game.
The real game returns a noise-flooded sample centered at
$Dec(sk,T_i.c)$.
The simulated-decrypt game returns a noise-flooded sample centered at the
stored plaintext $T_i.m_b$.
The formal proof factors the deterministic prefix of \textsc{ODec} from this
last noisy continuation, which is why the one-call replacement theorem can
focus only on the two Gaussian centers.

\begin{algorithm}
\caption{IND-CPA reduction adversary $\mathcal{B}_{\cA}$}
\label{alg:ind-cpa-reduction}
This is the conventional reduction endpoint used in the main proof.
It stores the same logical table as the IND-CPAD challenger, but obtains
challenge encryptions from the outer IND-CPA encryption oracle.
\begin{algorithmic}[1]
  \Procedure{$\mathsf{Flood}$}{$m,e$}
    \State $\eta \gets D_{\bbZ^n,(e^2+1)\gamma}$
    \State \Return $I_m^{-1}(\eta+I_m(m))$
  \EndProcedure
  \Statex
  \Procedure{$\mathcal{B}_{\cA}$}{$pk,evk$}
    \State $T \gets [\,]$
    \State $d \gets 0$
    \State $b^* \gets
      \cA^{OEnc_{\mathcal{B}},OEval_{\mathcal{B}},ODec_{\mathcal{B}}}(pk,evk)$
    \State \Return $b^*$
  \EndProcedure
  \Statex
  \Procedure{$OEnc_{\mathcal{B}}$}{$m_0,m_1$}
    \State $c \gets OEnc_b(m_0,m_1)$
      \Comment{outer IND-CPA encryption oracle}
    \State $T \gets T \doubleplus [(m_0,m_1,c)]$
    \State \Return $c$
  \EndProcedure
  \Statex
  \Procedure{$OEval_{\mathcal{B}}$}{$f,\mathcal{I}$}
    \State $\vec r \gets (T_i)_{i\in\mathcal{I}}$
    \State $m_0' \gets f((r.m_0)_{r\in\vec r})$
    \State $m_1' \gets f((r.m_1)_{r\in\vec r})$
    \State $c' \gets Eval(f,(r.c)_{r\in\vec r})$
    \State $T \gets T \doubleplus [(m_0',m_1',c')]$
    \State \Return $c'$
  \EndProcedure
  \Statex
  \Procedure{$ODec_{\mathcal{B}}$}{$i$}
    \State \textbf{assert} $d<q$
    \State $d \gets d+1$
    \State \textbf{assert} $i<|T|$
    \State $(m_0,m_1,c) \gets T_i$
    \If{$m_0\ne m_1$}
      \State \Return $\mathsf{None}$
    \EndIf
    \State \textbf{assert} $c=\mathsf{Some}(\_,e)$
    \State $y \gets \mathsf{Flood}(m_0,e)$
    \State \Return $\mathsf{Some}(y)$
  \EndProcedure
\end{algorithmic}
\end{algorithm}

The reduction does not know the IND-CPA challenge bit.
This is harmless for decryption queries that return a value: the simulator
only returns a flooded plaintext when the stored labels agree, $m_0=m_1$, so
there is no bit-dependent plaintext choice left to make.

\section{SSProve Syntax and Semantics}
\label[appendix]{app:ssprove-substrate}

This appendix recalls the fragment of SSProve on which the program logic is
built~\cite{ssprove}.
SSProve embeds pure expressions shallowly in Rocq and
represents effects by a free monad.
Fix a result type \(A\).
Its raw code has
the following constructors:
\begin{equation}
\label{eq:ssprove-syntax}
\begin{split}
c ::= {}& \mathsf{return}(a)
 \mid \mathsf{call}(p,x,\kappa)
 \mid \mathsf{get}(\ell,\kappa)\\
 &{}\mid \mathsf{put}(\ell,v,c)
 \mid \mathsf{sample}(D,\kappa).
\end{split}
\end{equation}
Here \(a:A\).
An operation signature \(p:S\to T\) contains a procedure
identifier as well as its argument and result types; a call has \(x:S\) and
\(\kappa:T\to\mathsf{RawCode}(A)\).
A typed location
\(\ell:T_\ell\) is read into a continuation
\(\kappa:T_\ell\to\mathsf{RawCode}(A)\), or updated with \(v:T_\ell\).
Finally, sampling takes \(D\in\mathsf{SD}(X)\) and
\(\kappa:X\to\mathsf{RawCode}(A)\), where \(\mathsf{SD}(X)\) denotes discrete
subdistributions on \(X\).
Continuations are Rocq functions, so conditionals,
pure local computation, and structurally terminating loops are supplied by
the ambient language rather than by additional constructors.
Monadic bind
recursively pushes a continuation through this syntax.
We use the customary
notations
\begin{equation}
\label{eq:ssprove-notation}
\begin{aligned}
 x\leftarrow c_1;c_2
   &\equiv\mathsf{bind}(c_1,\lambda x.c_2),\\
 x\leftarrow p(a);c
   &\equiv\mathsf{call}(p,a,\lambda x.c),\\
 x\leftarrow\mathsf{get}\,\ell;c
   &\equiv\mathsf{get}(\ell,\lambda x.c),\\
 \mathsf{put}\,\ell:=v;c
   &\equiv\mathsf{put}(\ell,v,c),\\
 x\leftarrow D;c
   &\equiv\mathsf{sample}(D,\lambda x.c).
\end{aligned}
\end{equation}

An interface is a finite set of typed operation signatures.
The checked type
\(\mathsf{Code}_{L,I}(A)\) pairs raw code with a proof that every accessed
location belongs to the footprint \(L\) and every external call belongs to
the import interface \(I\).
A package \(P:I\to E\) is a finite map from the
procedure names in its export interface \(E\) to well-scoped implementations
\(S\to\mathsf{Code}_{L,I}(T)\).
Sequential composition \(P\circ Q\)
links \(P\)'s calls to procedures implemented by \(Q\).
At code level, the
essential clause is
\begin{equation}
\label{eq:ssprove-link}
\begin{aligned}
 &\mathsf{codeLink}(\mathsf{call}(p,a,\kappa),Q)\\
 &\qquad =
 x\leftarrow Q.p(a);
 \mathsf{codeLink}(\kappa(x),Q),
\end{aligned}
\end{equation}
with \(\mathsf{codeLink}\) acting recursively on the remaining constructors.
Parallel composition combines disjoint export interfaces, while the identity
package forwards every call.
Package validity tracks the import, export, and
memory-footprint side conditions used by these operations.

Each location carries a type and a default initial value.
A heap \(h\) maps
locations to values of their declared types; lookup of an unset location
returns its default.
After linking has eliminated external calls, fully linked code denotes a
state-transforming subdistribution
\begin{equation}
\label{eq:ssprove-denotation}
 \semantic{c}:\mathsf{Heap}\longrightarrow
  \mathsf{SD}(A\times\mathsf{Heap}).
\end{equation}
Writing \(\delta_z\) for the point distribution at \(z\), its defining
equations are
\begin{align}
 \semantic{\mathsf{return}(a)}_h
   &=\delta_{(a,h)},\notag\\
 \semantic{\mathsf{get}(\ell,\kappa)}_h
   &=\semantic{\kappa(h(\ell))}_h,\notag\\
 \semantic{\mathsf{put}(\ell,v,c)}_h
   &=\semantic{c}_{h[\ell\mapsto v]},\notag\\
 \semantic{\mathsf{sample}(D,\kappa)}_h
   &=\sum_{x}D(x)\,\semantic{\kappa(x)}_h.
\label{eq:ssprove-effect-semantics}
\end{align}
The sums are pointwise sums of subdistributions.
In particular, semantic
bind threads both the returned value and the updated heap:
\begin{align}
 \semantic{x\leftarrow c;\kappa(x)}_h
 &=
 \sum_{(a,h')}
   \semantic{c}_h(a,h')\,
   \semantic{\kappa(a)}_{h'}.
\label{eq:ssprove-bind-semantics}
\end{align}
SSProve defines failure by sampling the null subdistribution, and
\(\mathsf{assert}(b)\) as \(\mathsf{return}(())\) when \(b\) holds and failure
otherwise.
Thus failed assertions have the zero subdistribution and their
continuations are not run.
A fully linked game, with no unresolved oracle calls, is evaluated by resolving
its \(\mathsf{Run}\) procedure, starting from the default heap, and projecting
the result component of \Cref{eq:ssprove-denotation}.

\input{appendix-checked-rules}
\section{Verified KL-Divergence Analysis}
\label{sec:verified-kl-analysis}

The program logic in the next section rests on a probability library rather than on informal analytic facts.
This section summarizes that library.
Its role is deliberately supporting: here we explain the mathematical facts that are available to the logic, while the next section explains how those facts are consumed by SSProve programs.

\subsection{KL Infrastructure}

The formalization defines KL divergence for MathComp distributions,
absolute continuity, and a finite-KL predicate used throughout the development:
\[
  \mathsf{finiteKL}(\cP,\cQ)
  \quad\Longleftrightarrow\quad
  \mathsf{absCont}(\cP,\cQ)
  \wedge
  \mathsf{summable}\!\left(
    x\mapsto
    \cP(x)\ln\!\left(\frac{\cP(x)}{\cQ(x)}\right)
  \right).
\]
The summability component matters in Rocq because many of the distributions in the development are countable rather than finite.
The core library also proves the log-sum inequality, KL preservation under injective pushforwards, and the summability facts used by the chain-rule proof.

The additive-error probability layer defines completion of a subdistribution by adding an explicit failure point,
constructs overlap and residual distributions, and proves a maximal-coupling theorem from a TVD bound:
\[
  \Delta(\mu,\nu)\leq \varepsilon
  \quad\Longrightarrow\quad
  \exists d.\ 
  \mathsf{Cpl}(d;\mu,\nu)
  \wedge
  \Pr_{(x,y)\gets d}[x=y]\geq 1-\varepsilon .
\]
The program-level additive-error rules use this theorem after packing SSProve outputs and heaps into a common completed output space.

\subsection{Pinsker and Pythagorean Preservation}

The formalization proves Pinsker's inequality for MathComp distributions:
\[
\begin{gathered}
  \mathsf{finiteKL}(\cP,\cQ)
  \wedge
  dweight(\cP)=dweight(\cQ)=1
\\
  \Longrightarrow
  \Delta(\cP,\cQ)
  \leq
  \sqrt{\frac{\KL{\cP}{\cQ}}{2}}.
\end{gathered}
\]
This is not assumed axiomatically.
The proof reduces total variation to a binary event, proves the binary quadratic Pinsker inequality by real analysis, and combines it with the data-processing/log-sum argument for that event.

The chain-rule development proves the conditional-coordinate KL bound.
This is the result used by Pythagorean probability preservation.
For distributions $\cP,\cQ$ on $n$-tuples, the theorem has the following shape:
\[
\begin{gathered}
  dweight(\cP)=dweight(\cQ)=1,
  \qquad
  \forall i,a.\ 
    \mathsf{finiteKL}(\cP_i\mid a,\cQ_i\mid a),
\\
  \forall i,a.\ 
    \KL{\cP_i\mid a}{\cQ_i\mid a}\leq \varepsilon_i
  \quad\Longrightarrow\quad
  \KL{\cP}{\cQ}\leq \sum_i \varepsilon_i .
\end{gathered}
\]
The proof is the technical heart of the probability development:
it decomposes the KL integrand pointwise over tuple prefixes, controls both positive and negative parts of the resulting series, and then reassembles the bound.
The value-level inequality and the summability obligation are proved in parallel.
Both use the same grouping step: for each coordinate, traces are grouped by their prefix and current coordinate value, while the suffix is summed away.
For the value theorem this produces an expectation of conditional KL divergences; for summability the same grouping is applied to finite positive-part sums, with one negative-part slack term per coordinate.

Combining this chain bound with Pinsker gives the main preservation theorem.
\begin{lemma}[Pythagorean probability preservation]
\label{lem:verified-pythagorean-preservation}
Let $\cP,\cQ\in\distr(A^n)$ be distributions of weight one, and let
$s=(s_1,\ldots,s_n)$ be a tuple of nonnegative real numbers.
If, for every coordinate $i$ and every prefix $a\in A^{i-1}$,
\[
  \mathsf{finiteKL}(\cP_i\mid a,\cQ_i\mid a)
  \qquad\text{and}\qquad
  \KL{\cP_i\mid a}{\cQ_i\mid a}\leq s_i,
\]
then
\[
  \Delta(\cP,\cQ)
  \leq
  \sqrt{\frac{\sum_i s_i}{2}}.
\]
\end{lemma}
The same part of the formalization packages this lemma into a reusable Pythagorean distribution predicate, proves reflexivity and singleton-trace lemmas, and defines the call-error tuples used by the trace-compiler theorem.

\subsection{Discrete Gaussian Facts}

The discrete-Gaussian development is now one ingredient of the KL story rather than the whole story.
It proves the equal-variance KL formula used by \LabTirName{KL-DG}:
\[
  \sigma>0
  \quad\Longrightarrow\quad
  \KL{\DG{\mu}{\sigma^2}}{\DG{\nu}{\sigma^2}}
  =
  \frac{(\nu-\mu)^2}{2\sigma^2}.
\]
The proof is completely internal to the development.
It first exposes the normalized mass function, proves integer-shift invariance of the Gaussian weight, derives symmetry of the centered distribution, and then uses the zero-centered first moment to evaluate the remaining expectation.
The moment summability needed for that expectation is proved as well; it is no longer a side caveat.

The development also includes geometric upper bounds for events of the form
\[
  \Pr_{x\gets\DG{\mu}{\sigma^2}}
    \bigl[\, |x-\mu|\geq k\,\bigr].
\]
These tail bounds are useful for future cryptographic applications, but the current Pythagorean program logic primarily consumes the KL formula above through the discrete-Gaussian sampling rule.

Thus the probability layer supplies three kinds of facts to the program logic:
ordinary distribution-level KL bounds such as the discrete-Gaussian theorem,
global KL-to-TVD conversion via Pinsker and the chain rule,
and completion/coupling lemmas that turn TVD bounds into additive-error judgments over SSProve semantics.

\section{Formal Interface Excerpts}
\label{app:formal-interface-excerpts}

This appendix records selected Rocq excerpts that determine the meaning of the
main theorem.
The excerpts are included from a paper-local snapshot of the artifact source
files, refreshed by the paper build.
This keeps the paper self-contained for Overleaf while preserving a simple
path back to the canonical Rocq development.
They are not intended to replace the full artifact; rather, they expose the
small specification surface that a cryptographic reader should be able to
audit independently.

\begingroup
\setmonofont{FreeMono}[Scale=MatchLowercase]

\subsection{Approximate Homomorphic Encryption Scheme Interface}
\label{app:approximate-homomorphic-encryption-interface}

\inputminted[
  fontsize=\scriptsize,
  breaklines,
  linenos,
  firstnumber=25,
  firstline=25,
  lastline=50
]{coq}{rocq-excerpts/theories/Schemes/ApproxFHE.v}

\subsection{Metric Charts}
\label{app:metric-chart-excerpt}

\inputminted[
  fontsize=\scriptsize,
  breaklines,
  linenos,
  firstnumber=105,
  firstline=105,
  lastline=121
]{coq}{rocq-excerpts/theories/Schemes/ApproxFHE.v}

\subsection{Support-Level Correctness}
\label{app:support-correctness-excerpt}

\inputminted[
  fontsize=\scriptsize,
  breaklines,
  linenos,
  firstnumber=191,
  firstline=191,
  lastline=230
]{coq}{rocq-excerpts/theories/Schemes/ApproxFHE.v}

\subsection{IND-CPA Game}
\label{app:ind-cpa-game-excerpt}

\inputminted[
  fontsize=\scriptsize,
  breaklines,
  linenos,
  firstnumber=53,
  firstline=53,
  lastline=112
]{coq}{rocq-excerpts/theories/Schemes/Indcpa.v}

\subsection{IND-CPA with Decryption Oracle and Game}
\label{app:ind-cpa-d-game-excerpt}

\inputminted[
  fontsize=\scriptsize,
  breaklines,
  linenos,
  firstnumber=66,
  firstline=66,
  lastline=130
]{coq}{rocq-excerpts/theories/Schemes/Indcpad.v}

\inputminted[
  fontsize=\scriptsize,
  breaklines,
  linenos,
  firstnumber=138,
  firstline=138,
  lastline=180
]{coq}{rocq-excerpts/theories/Schemes/Indcpad.v}

\subsection{Noise-Flooded Construction}
\label{app:noise-flooded-construction-excerpt}

\inputminted[
  fontsize=\scriptsize,
  breaklines,
  linenos,
  firstnumber=258,
  firstline=258,
  lastline=302
]{coq}{rocq-excerpts/theories/Constructions/NoiseFlooding.v}

\subsection{Concrete IND-CPA Reduction}
\label{app:ind-cpa-reduction-excerpt}

\inputminted[
  fontsize=\scriptsize,
  breaklines,
  linenos,
  firstnumber=224,
  firstline=224,
  lastline=318
]{coq}{rocq-excerpts/theories/Security/IndcpadSimulator.v}

\subsection{Loss Formula}
\label{app:loss-formula-excerpt}

\inputminted[
  fontsize=\scriptsize,
  breaklines,
  linenos,
  firstnumber=39,
  firstline=39,
  lastline=43
]{coq}{rocq-excerpts/theories/Security/NoiseFloodingSecurity/Prelude.v}

\subsection{Security Bound}
\label{app:security-bound-excerpt}

\inputminted[
  fontsize=\scriptsize,
  breaklines,
  linenos,
  firstnumber=546,
  firstline=546,
  lastline=549
]{coq}{rocq-excerpts/theories/Security/NoiseFloodingSecurity/Final.v}

\subsection{Reduction Validity}
\label{app:reduction-bound-excerpt}

\inputminted[
  fontsize=\scriptsize,
  breaklines,
  linenos,
  firstnumber=2375,
  firstline=2375,
  lastline=2379
]{coq}{rocq-excerpts/theories/Security/NoiseFloodingSecurity/Final.v}

\subsection{Final Security Theorem}
\label{app:final-reduction-theorem}

\inputminted[
  fontsize=\scriptsize,
  breaklines,
  linenos,
  firstnumber=2607,
  firstline=2607,
  lastline=2610
]{coq}{rocq-excerpts/theories/Security/NoiseFloodingSecurity/Final.v}

\endgroup


\end{document}
